\chapter{Surgery}

\medterm{Acute Mesenteric Ischemia}-- Sudden interruption of blood flow to the small intestine, with mortality 60-80\% if diagnosis delayed. Etiologies: arterial embolism (50\%, SMA most common, atrial fibrillation, recent MI, valvular disease), arterial thrombosis (25\%, on pre-existing atherosclerosis, often at SMA origin), venous thrombosis (20\%, hypercoagulable states, malignancy, OCPs, pregnancy), non-occlusive mesenteric ischemia (NOMI, 5\%, low flow state: shock, heart failure, vasopressors). prognosis, with second-look laparotomy to reassess intestinal viability. If extensive infarction, the survival rate is dismal, and palliative care should be considered.

\medterm{Acute Pancreatitis}-- Acute pancreatitis (AP): autodigestion of pancreas by prematurely activated trypsin. Causes: gallstones (40-50\%), alcohol (25-35\%), hypertriglyceridemia (>1000 mg/dL), ERCP, medications (azathioprine, 6-MP, valproate, estrogens, thiazides), trauma, post-surgical, autoimmune, idiopathic. Severity: Atlanta classification 2012: mild (no organ failure, no local complications), moderately severe (transient organ failure <48h or local complications), severe (persistent organ failure >48h). Revised Atlanta classification includes three severity categories. Local complications include acute peripancreatic fluid collections (APFC), pancreatic pseudocysts (developing >4 weeks after onset), acute necrotic collections (ANC), and walled-off necrosis (WON). Pancreatic and peripancreatic necrosis may become infected in 30-70\% of cases after 3-4 weeks, requiring antibiotics and percutaneous or endoscopic drainage. Management: aggressive IV fluids (lactated Ringer's, 250-500 mL/h for first 12-24 hours), analgesia, NPO with NG tube for vomiting, ICU monitoring for severe disease. Antibiotics are not routinely indicated. ERCP with sphincterotomy is indicated for severe gallstone pancreatitis with cholangitis or persistent choledocholithiasis. Necrosectomy (minimally invasive step-up approach via percutaneous drainage, endoscopic transluminal drainage, video-assisted retroperitoneal debridement, or surgical necrosectomy) is reserved for infected necrosis not responding to antibiotics. Complications include multi-organ failure, ARDS, acute kidney injury, splenic vein thrombosis, pseudoaneurysm formation with hemorrhage, and chronic pancreatitis with exocrine and endocrine insufficiency. Prognosis: CT severity index (CTSI), APACHE II, Ranson criteria, and BISAP score predict severity. Mortality of severe acute pancreatitis with persistent organ failure is 30-50\%.

\medterm{Adrenal Tumors}-- Adrenal incidentaloma: >1cm on imaging. Workup: hormonal (cortisol after 1mg DST, metanephrines, aldosterone/renin if hypertensive, DHEAS/androgens if virilizing), imaging (CT washout: adenoma >60\% absolute washout or >40\% relative; MRI chemical shift). Size >4cm or suspicious imaging -> surgery. Pheochromocytoma: catecholamine excess (episodic hypertension, headache, diaphoresis, palpitations); 10\% bilateral, 10\% malignant, 10\% extra-adrenal (paraganglioma), 10\% familial (MEN2, VHL, NF1, SDHx). Biochemical confirmation: plasma free metanephrines (elevated). Pre-operative: alpha-blockade (phenoxybenzamine 10 mg TDS for 7-14 days before surgery, starting at 10 mg TDS and titrating up to 20-30 mg TDS to achieve blood pressure control and prevent intraoperative hypertensive crisis, followed by beta-blockade after alpha-blockade to control reflex tachycardia or arrhythmia, such as propranolol 10-20 mg TDS starting on day 3 of alpha-blockade). Surgery: laparoscopic adrenalectomy (gold standard for pheochromocytoma and other adrenal masses; partial/cortical-sparing adrenalectomy if bilateral disease or genetic syndrome). Cortical-sparing adrenalectomy is preferred for hereditary phaeochromocytoma syndromes (MEN2, VHL) to avoid lifelong steroid dependence. Criteria for surgery: tumors >4 cm or any functioning or suspicious mass in a patient who is a good surgical candidate. Adrenocortical carcinoma (ACC): usually >6cm, irregular, heterogeneous with central necrosis, invasion of adjacent structures, venous tumor thrombus; worst prognosis; treatment is open radical en bloc resection with lymphadenectomy and complete tumor excision achieving R0 resection margins. Mitotane therapy for advanced or metastatic ACC. Adjuvant mitotane therapy after resection reduces recurrence when there is a high risk of recurrence (tumor size >6 cm, Ki-67 >10 percent, positive margins, capsular invasion, or tumor rupture at the time of surgery).

\medterm{Amputation}-- Surgical removal of a limb or part of a limb. Levels: toe, partial foot (transmetatarsal), below-knee (BKA, preserves knee function, 40-60\% walk with prosthesis), through-knee, above-knee (AKA), hip disarticulation, and hemipelvectomy. Upper extremity: digit, ray, transradial, transhumeral, and shoulder disarticulation. Indications: peripheral arterial disease with gangrene (most common), trauma (crush, degloving), infection (necrotizing fasciitis, osteomyelitis refractory), malignancy (sarcoma), and severe deformity. Prosthetic fitting: post-operative rigid dressing, early weight-bearing for BKA. Complications: stump pain, phantom limb pain (80-100\%), phantom sensation, surgical site infection, and deconditioning.

\medterm{Anti-Reflux Surgery}-- Surgical procedures designed to restore the anti-reflux barrier at the gastroesophageal junction, primarily indicated for gastroesophageal reflux disease (GERD) refractory to medical therapy, patient preference to discontinue lifelong medication, or large hiatal hernias. The most common procedure is laparoscopic Nissen fundoplication, a 360-degree wrap of the gastric fundus around the distal esophagus, restoring the lower esophageal sphincter pressure and preventing transient lower esophageal sphincter relaxation and reducing the volume of refluxate. The wrap is typically 2-3 cm in length and is performed laparoscopically. Pre-operative evaluation includes esophageal manometry to confirm normal peristalsis (essential before a 360-degree wrap), 24-hour pH monitoring to confirm the diagnosis of GERD, upper endoscopy to assess for oesophagitis, Barrett esophagus, and stricture, and a barium swallow to evaluate esophageal anatomy and identify hiatal hernia. Alternatives to Nissen fundoplication include partial wraps (Toupet 270-degree posterior fundoplication for patients with poor esophageal motility, Dor 180-200-degree anterior fundoplication) and the LINX magnetic sphincter augmentation device, which consists of a ring of titanium beads with magnetic cores placed around the gastro-esophageal junction. Crural closure of the diaphragmatic hiatus is performed to prevent recurrence of the hiatal hernia. Post-operative complications include dysphagia (common initially, resolves in most cases, may require pneumatic dilation if persistent), gas-bloat syndrome (inability to belch, leading to abdominal distension and discomfort, usually improves over weeks to months), and recurrence of reflux symptoms due to wrap disruption or slippage. Outcomes are excellent, with 85-95 percent patient satisfaction at 10 years.

\medterm{Appendectomy}-- Surgical removal of the appendix, the most common emergency surgical procedure performed worldwide. Appendectomy is indicated for acute appendicitis, the most common surgical emergency with a lifetime risk of 7-8 percent and peak incidence between 10 and 30 years. Laparoscopic appendectomy is the preferred approach, offering lower wound infection rates, faster recovery, less postoperative pain, and better cosmetic outcomes compared to open appendectomy. The procedure involves placement of three to four ports (umbilical camera port, suprapubic and right lower quadrant working ports), identification of the base of the appendix, division of the mesoappendix using cautery or a vessel-sealing device, and ligation of the appendiceal base using endoloops (two proximal and one distal) or a stapler. The appendix is placed in a retrieval bag and removed through the umbilical port. In cases of perforated appendicitis with peritonitis, conversion to open appendectomy may be required for adequate peritoneal lavage. Intravenous antibiotics covering Gram-negative and anaerobic organisms are given pre-operatively. For uncomplicated appendicitis, no post-operative antibiotics are needed beyond 24 hours. For perforated appendicitis, antibiotics are continued for 3-5 days or until the patient is clinically improving. Complicated appendicitis (perforation, abscess, or gangrene) has a higher risk of wound infection and intra-abdominal abscess, which may require percutaneous drainage. The complication rate for laparoscopic appendectomy is approximately 5-10 percent, with wound infection being the most common (1-5 percent). Post-operative recovery is rapid, with most patients discharged within 24 hours and returning to normal activity within 1-2 weeks.

\medterm{Appendiceal Tumors}-- Appendiceal neoplasms: carcinoid (most common, 50\%, tip of appendix, <1cm: appendectomy alone; >2cm or positive margins/lymphovascular invasion: right hemicolectomy), mucinous neoplasm (low-grade appendiceal mucinous neoplasm LAMN, high-grade appendiceal mucinous neoplasm HAMN), appendiceal adenocarcinoma (treated like colon cancer: right hemicolectomy), goblet cell carcinoid (adenocarcinoid, amphicrine, right hemicolectomy if >2cm or nodal involvement). Pseudomyxoma peritonei (PMP): disseminated mucinous tumor throughout the peritoneal cavity from ruptured mucinous neoplasm; treatment: cytoreductive surgery (CRS) + hyperthermic intraperitoneal chemotherapy (HIPEC) with mitomycin C. Prognosis: 5-year survival 70\% for low-grade, 20-50\% for high-grade.

\medterm{Appendicitis}-- Acute appendicitis is the most common surgical emergency, with a lifetime risk of 7-8 percent. Peak incidence 10-30 years. Pathophysiology: obstruction of appendiceal lumen (fecalith, lymphoid hyperplasia, tumor, parasites) -> increased intraluminal pressure -> venous congestion -> arterial compromise -> ischemia, perforation, peritonitis. Clinical: periumbilical pain migrating to RLQ (McBurney point), anorexia, nausea/vomiting, low-grade fever. Signs: RLQ tenderness, rebound tenderness, Rovsing sign (palpation LLQ causes pain in RLQ), psoas sign (pain on right hip extension and thigh flexion), obturator sign (pain on internal rotation of flexed right thigh), Dunphy sign (cough worsens RLQ pain), and Aaron sign (referred pain in epigastrium or precordium when palpating RLQ). Laboratory: leukocytosis (WBC >10,000) with left shift, but WBC may be normal in 10-20\% of cases; C-reactive protein (CRP) elevated; procalcitonin may be used. Imaging: graded compression ultrasound (first-line, sensitivity 86\%, specificity 93\%) shows non-compressible, blind-ending tubular structure >6 mm in diameter, with increased blood flow on color Doppler. CT abdomen with IV contrast is the gold standard for diagnosis (sensitivity 94\%, specificity 95\%), showing an enlarged appendix >6 mm with wall thickening, periappendiceal fat stranding, and possible abscess or extraluminal air (indicating perforation). Alternative diagnoses to exclude include mesenteric adenitis, terminal ileitis (Crohn disease), Meckel diverticulitis, right-sided diverticulitis, omental infarction, epiploic appendagitis, ureteric colic, pyelonephritis, ovarian cyst torsion or rupture, ectopic pregnancy, pelvic inflammatory disease, and mitelschmerz. The Alvarado score (MANTRELS) and RIPASA score stratify risk. Alvarado score: 1-4 = low risk (observe), 5-6 = moderate (CT), 7-10 = high (surgery). Management: appendectomy (laparoscopic preferred). Antibiotics alone (appendicitis non-operative management) is an option for uncomplicated (non-perforated, no abscess, no faecalith) appendicitis based on the NOTA, APPAC, and CODA trials, but 20-30\% recurrence rate within 1 year. Interval appendectomy after non-operative management is no longer routinely recommended. Complications: perforation (incidence 20-30\%, higher in elderly and children), abscess formation (percutaneous drainage), wound infection (1-5\%), intra-abdominal abscess (1-2\%), ileus, and stump leak (rare).

\medterm{Arterial Line (A-Line)}-- A catheter inserted into a peripheral artery (most commonly radial, also femoral, brachial, dorsalis pedis) for continuous blood pressure monitoring and arterial blood gas sampling. Allen test: assesses ulnar collateral flow before radial artery cannulation. Indications: hemodynamic instability (shock, vasopressors), need for frequent ABGs, and advanced hemodynamic monitoring (pulse pressure variation SVV for fluid responsiveness). Transducer: zeroed at the phlebostatic axis (4th intercostal space, midaxillary line), leveled to the patient's heart. Damping: underdamped (overestimates systolic, normal waveform) or overdamped (underestimates systolic, sluggish upstroke). Complications: thrombosis, ischemia (rare if collateral flow is adequate), infection, hematoma, and pseudoaneurysm.

\medterm{Arterial Ulcer}-- A chronic wound caused by arterial insufficiency (peripheral artery disease). Located on distal extremities: toes, heels, lateral malleolus, and sites of pressure. Presents with punched-out, deep, well-demarcated wound with pale or necrotic base, minimal exudate, and cool, hairless, shiny skin. Pain is severe, worse with elevation and night. Associated with claudication and rest pain. Diagnosis: ABI (<0.5 suggests severe ischemia), toe pressure (<30 mmHg), vascular duplex. Treatment: revascularization (endovascular or surgical bypass) is essential for healing. Wound care: debridement, moist healing, infection control. Avoid compression. Amputation if non-salvageable.

\medterm{Bariatric Surgery}-- Weight loss procedures and surgery, collectively known as bariatric surgery, are invasive interventions for severe obesity that produce substantial and sustained weight loss by altering gastrointestinal anatomy and physiology. Common procedures include Roux-en-Y gastric bypass, sleeve gastrectomy, and adjustable gastric banding, which restrict food intake, reduce nutrient absorption, or both. Bariatric surgery is typically reserved for individuals with a body mass index of 40 or greater, or 35 or greater with significant obesity-related complications, and requires lifelong nutritional monitoring and lifestyle adherence.

\medterm{Benign Prostatic Hyperplasia}-- Non-malignant enlargement of the prostate gland due to hyperplasia of stromal and glandular elements, affecting >50\% of men by age 60 and >90\% by age 85. Presents with lower urinary tract symptoms (LUTS): hesitancy, weak stream, intermittency, straining, incomplete emptying, urgency, frequency, and nocturia. Complications: acute urinary retention, bladder stones, recurrent UTIs, and renal impairment. Diagnosis: digital rectal exam, PSA, AUA symptom index, post-void residual, and uroflowmetry. Treatment: watchful waiting for mild symptoms, alpha-blockers (tamsulosin), 5-alpha-reductase inhibitors (finasteride), combination therapy, and surgical options (TURP, laser prostatectomy, and open prostatectomy for large glands).

\medterm{Biliary Tract Injuries}-- Bile duct injury (BDI) occurs in 0.3-0.5\% of laparoscopic cholecystectomies. Strasberg classification: A (cystic duct leak/minor hepatic duct leak), B (occlusion of aberrant right hepatic duct), C (transection of aberrant right hepatic duct without ligation), D (lateral injury to major duct), E (major duct injury: E1-E5 based on level). Recognition: intraoperative (bile leak, visualization) or postoperative (bilious drain output, pain, fever, jaundice). Intraoperative management: if recognized, primary repair over T-tube (if >50\% intact wall) or Roux-en-Y hepaticojejunostomy (if transection, ligation, or significant resection of the bile duct). If recognized postoperatively, ERCP with sphincterotomy and stenting (Strasberg A, D) or percutaneous transhepatic cholangiography (PTC) with drainage (Strasberg B, C) may be attempted. Strasberg E injuries require surgical repair via Roux-en-Y hepaticojejunostomy, preferably performed by a hepatobiliary surgeon at a tertiary referral center. The timing of repair is critical: early repair (<72 hours) is acceptable for clean, well-vascularised tissue; otherwise, delayed repair (4-6 weeks) after resolution of inflammation and sepsis is preferred. Pre-operative optimisation includes biliary drainage (ERCP, PTC, or percutaneous transhepatic biliary drainage), control of sepsis with antibiotics (piperacillin-tazobactam or meropenem), and nutritional support (enteral nutrition via nasojejunal tube if possible). Prevention: critical view of safety (CVS) technique during cholecystectomy, low threshold for cholangiography (indocyanine green fluorescence imaging, intraoperative cholangiography, or laparoscopic ultrasound) when anatomy is unclear. Outcomes: Strasberg A and D have excellent outcomes; Strasberg E has a 10-20\% stricture rate after repair. Long-term complications include anastomotic stricture (treated with balloon dilation or revision), cholangitis, secondary biliary cirrhosis, and portal hypertension.

\medterm{Blunt Abdominal Trauma}-- Trauma to the abdomen without penetrating injury. Mechanisms: motor vehicle collisions (most common), falls, assaults. Commonly injured organs: spleen (most common, left upper quadrant), liver (right upper quadrant), kidneys (flank), bowel (seatbelt sign, delayed perforation), and pancreas (epigastric). Diagnosis: FAST (Focused Assessment with Sonography in Trauma) for free intraperitoneal fluid; CT scan (gold standard in hemodynamically stable patients). Management: hemodynamically unstable with positive FAST → exploratory laparotomy. Stable: CT for grading; non-operative management of solid organ injuries (spleen/liver) if hemodynamically stable. Bowel and pancreatic injuries often require surgical repair.

\medterm{BPH}-- Abbreviation for Benign Prostatic Hyperplasia.

\medterm{Breast Surgery}-- Breast cancer surgery: breast-conserving therapy (BCT: lumpectomy + radiation) equivalent survival to mastectomy for early stage. Margins: "no ink on tumor" for invasive cancer; 2mm for DCIS. Sentinel lymph node biopsy (SLNB): dual mapping (blue dye + radioisotope); if positive (macrometastasis >2mm): completion ALND (historical) or radiation (ACOSOG Z0011: 1-2 positive SLN + BCT + whole breast RT -> no ALND; AMAROS: ALND vs axillary RT equivalent). Mastectomy types: simple, skin-sparing, nipple-sparing (NSM), and radical modified (Patey: breast + axillary lymph node dissection - ALND). Nipple-areola complex (NAC) is spared if tumor is >2cm from the nipple and no Paget disease of the breast. Oncoplastic techniques integrate plastic surgery principles with oncological resection for improved aesthetic outcomes. Adjuvant therapy: radiation (post-BCT, post-mastectomy if T3/T4 or 4+ positive nodes), endocrine (tamoxifen pre-menopausal, aromatase inhibitor post-menopausal if ER+), HER2-targeted (trastuzumab/pertuzumab if HER2+), chemotherapy (high-risk: triple-negative, HER2+, high-grade, node-positive). BRCA mutation carriers: bilateral mastectomy (risk-reducing) + RRSO (risk-reducing salpingo-oophorectomy) recommended.

\medterm{Burn Surgery}-- Burn assessment: rule of 9s (adult), Lund-Browder (children). Depth: superficial (epidermis), superficial partial (dermis, blanch, blister, painful), deep partial (deeper dermis, white, sluggish capillary refill, less painful), full thickness (all dermis, white/charred, anesthetic). TBSA >20\%: IV resuscitation (Parkland: 4 mL/kg/\%TBSA LR, half in first 8h, half in next 16h; modified Brooke: 2 mL/kg/\%TBSA). Inhalation injury: bronchoscopy, carboxyhemoglobin, cyanide level. Escharotomy: circumferential full-thickness burns of chest (impaired ventilation), extremities (compartment syndrome), neck, or penis. Fasciotomy for compartment syndrome (burn deep to fascia). Tangential excision (within 24-72 hours) with split-thickness skin grafting (STSG: 0.008-0.012 inch), meshed (1:1 to 1:4 for larger surface area) or sheet graft (face, hands, joints). Biological dressings: allograft (cadaver), xenograft (porcine), amniotic membrane, synthetic (Biobrane, Integra, Apligraf, and artificial dermal regeneration template). Scar management: pressure garments, silicone sheeting, laser therapy (pulsed dye laser for hypertrophic scars, fractional CO2 laser for contracture release and scar texture improvement).

\medterm{CABG (Coronary Artery Bypass Grafting)}-- Open heart surgery to bypass stenotic coronary arteries using autologous grafts. On-pump: cardiopulmonary bypass, cardioplegic arrest. Off-pump (OPCAB): beating heart, no bypass machine. Grafts: left internal mammary artery (LIMA to LAD, best patency 95\% at 10 years), saphenous vein graft (SVG, patency 60-80\% at 10 years), and radial artery (patency between LIMA and SVG). Indications: left main CAD (\(\ge\)50\%), three-vessel CAD, two-vessel CAD with proximal LAD and EF <50\%, and failed PCI. Complications: bleeding, AF, stroke, infection (sternal wound), renal failure, and graft failure.

\medterm{Carotid Endarterectomy}-- Surgical removal of atherosclerotic plaque from the carotid bifurcation to prevent stroke. Indications: symptomatic carotid stenosis \(\ge\)50\% (NASCET criteria: TIA or non-disabling stroke) and asymptomatic stenosis \(\ge\)70\% (if perioperative risk <3\%). Pre-operative: CTA or MRA, antiplatelet therapy. Complications: stroke (2-5\%), cranial nerve injury (hypoglossal, vagus, facial marginal mandibular, recurrent laryngeal: 5-10\%, most temporary), hematoma, restenosis. Alternative: carotid artery stenting (CAS, for high surgical risk).

\medterm{Cesarean Section (C-section)}-- Surgical delivery of a fetus through incisions in the abdominal wall (laparotomy) and uterus (hysterotomy). Indications: failure to progress, non-reassuring fetal heart tracing, malpresentation (breech, transverse), placenta previa, placental abruption, prior classical C-section, maternal infection (HSV, HIV), and elective (maternal request). Incisions: low transverse (Kerr, most common, heals best for future VBAC), classical (vertical, high risk of rupture in labor). Steps: Pfannenstiel incision, bladder flap, hysterotomy, delivery of the baby, placental removal, closure. Complications: hemorrhage, infection, thromboembolism, bladder/bowel injury, and uterine rupture in future pregnancies (VBAC).

\medterm{Cholecystectomy}-- Surgical removal of the gallbladder. Laparoscopic cholecystectomy (gold standard, >95\%): 4 small incisions, short stay. Open cholecystectomy: for severe inflammation, cirrhosis, or failed laparoscopy. Indications: symptomatic cholelithiasis (biliary colic), acute cholecystitis, gallstone pancreatitis, and gallstone ileus. Intraoperative cholangiogram (IOC): assess for CBD stones. Critical view of safety (CVS): before clipping the cystic artery and duct. Complications: bile duct injury (0.3-0.5\%), bile leak, bleeding, and retained CBD stones.

\medterm{Cholecystitis}-- Acute cholecystitis is gallbladder inflammation, most commonly from cystic duct obstruction by gallstones (calculous cholecystitis, 90-95\%); acalculous cholecystitis (5-10\%) in critically ill, trauma, burns, prolonged fasting, TPN, vasculitis. Pathophysiology: stone impaction -> distension -> ischemia -> inflammation -> bacterial translocation (E. coli, Klebsiella, Enterococcus). Clinical: RUQ pain >6 hours, radiating to right scapula, fever, nausea/vomiting. Murphy sign (inspiratory arrest on right subcostal palpation during deep inspiration), present in 60-70\% of cases. Laboratory: leukocytosis, elevated CRP, LFTs (elevated ALT/AST in acute cholecystitis, ALP/bilirubin in choledocholithiasis or Mirizzi syndrome). Imaging: RUQ ultrasound (first-line, sensitivity 94\%): gallstones, thickened gallbladder wall >3-4 mm (target sign), pericholecystic fluid, sonographic Murphy sign, positive common bile duct >6 mm (or >8 mm post-cholecystectomy). CT abdomen is less sensitive for gallstones but useful for complications (perforation, abscess, emphysematous cholecystitis, gangrenous cholecystitis, Mirizzi syndrome). HIDA scan (cholescintigraphy) shows non-visualisation of the gallbladder (cystic duct obstruction) with 95\% sensitivity and specificity. Management: NPO, IV fluids, analgesia, IV antibiotics (ampicillin-sulbactam, piperacillin-tazobactam, or ceftriaxone with metronidazole). Definitive: laparoscopic cholecystectomy (gold standard), ideally within 72 hours of admission (early cholecystectomy reduces morbidity and length of stay compared to delayed surgery at 6-12 weeks). Tokyo Guidelines 2018 (TG18) severity grading: Grade I (mild, no organ dysfunction, mild inflammation), Grade II (moderate, elevated WBC >18,000, palpable tender mass, duration >72 hours, marked inflammation), Grade III (severe, organ dysfunction). Grade I: early laparoscopic cholecystectomy (within 72 hours); Grade II: early laparoscopic cholecystectomy after medical optimisation; Grade III: ICU management, percutaneous cholecystostomy (drainage) if critical, delayed cholecystectomy after resolution. Complications: gallbladder perforation (localised or free), empyema of gallbladder, emphysematous cholecystitis (gas-forming organisms, high mortality, requires urgent cholecystectomy), Mirizzi syndrome (stone impacted in cystic duct or Hartmann pouch compressing common hepatic duct causing obstructive jaundice; may require subtotal cholecystectomy or Roux-en-Y hepaticojejunostomy), post-cholecystectomy syndrome (persistent symptoms after cholecystectomy).

\medterm{Chronic Pancreatitis}-- Chronic pancreatitis (CP): irreversible fibrosis, loss of exocrine/endocrine function. Etiology: alcohol (60-70\%), genetic (PRSS1, SPINK1, CFTR, CTRC), obstructive (stricture, tumor), autoimmune (IgG4-related), idiopathic. Clinical: recurrent epigastric pain radiating to back, steatorrhea (exocrine insufficiency), diabetes (endocrine insufficiency), weight loss. Imaging: CT (calcifications, ductal dilation, atrophy), MRCP (ductal changes), EUS (early changes: hyperechoic foci, strands, lobularity), ERCP (therapeutic: stenting, stone extraction, stricture dilation). Cambridge classification (ERCP/MRCP) grade 1-5. Management: pain control (avoid narcotics, consider celiac plexus block for refractory pain), exocrine supplementation (pancreatic enzymes, PERT: lipase 40,000-80,000 units per meal, titrate to symptom relief and normalisation of fecal elastase-1), endocrine (insulin for diabetes), endoscopic (ERCP with sphincterotomy, stone extraction, stricture dilation, stenting), surgical (lateral pancreaticojejunostomy/Puestow procedure for dilated duct >7 mm, Beger/Frey procedure for head-predominant disease, Whipple for mass/head, distal pancreatectomy for body/tail disease, total pancreatectomy with islet autotransplantation for debilitating pain). Complications: pseudocyst (can be internal or external, may require endoscopic drainage, transpapillary drainage, or cystogastrostomy), biliary stricture, duodenal stricture, splenic vein thrombosis leading to left-sided (sinistral) portal hypertension with gastric varices, pancreatic ascites, pancreaticopleural fistula, and pancreatic adenocarcinoma (risk is 2-5\% over 20 years, surveillance of the remnant gland with cross-sectional imaging is controversial but may be considered in young patients).

\medterm{Circumcision}-- Surgical removal of the penile foreskin (prepuce). Newborn: most common, Plastibell or Gomco clamp, or Mogen clamp (bloodless). Indications: religious (Jewish Brit Milah), cultural, medical (phimosis, recurrent balanitis, balanitis xerotica obliterans, paraphimosis, and UTI prevention in boys with VUR). Medical benefits: reduced UTI risk (10x decrease in infancy), reduced penile cancer risk, lower HIV/STI transmission. Contraindications: hypospadias, epispadias, chordee (foreskin needed for repair). Complications: bleeding, infection, meatal stenosis, excess/insufficient removal, and urethrocutaneous fistula.

\medterm{Cleft Lip and Palate}-- A congenital anomaly caused by failure of fusion of the facial prominences during embryonic development (4-10 weeks gestation). Cleft lip: failure of maxillary and medial nasal processes to fuse. Cleft palate: failure of palatal shelves to fuse. Can occur together or separately. Incidence: 1 in 700 live births, more common in Asians and Native Americans. Multifactorial: genetic (IRF6, FGFR1) and environmental (maternal smoking, folate deficiency, alcohol). Associated syndromes: Pierre Robin sequence, Van der Woude syndrome, 22q11 deletion. Presents with cosmetic deformity, feeding difficulties, speech abnormalities, and ear infections. Management: multidisciplinary team (ENT, plastic surgery, speech therapy, orthodontics). Surgical repair: cleft lip at 3-6 months, cleft palate at 9-18 months.

\medterm{Cleft Palate}-- A congenital malformation characterized by an opening (cleft) in the palate (roof of the mouth) due to failure of fusion of the palatal shelves during embryonic development (weeks 6--9 of gestation). Incidence: 1 in 1,500--2,500 live births. Cleft palate may occur in isolation or with cleft lip (cleft lip and palate—CLP). Classification: incomplete (soft palate only) vs complete (soft and hard palate), unilateral vs bilateral, and associated with Pierre Robin sequence (micrognathia, glossoptosis, airway obstruction). Risk factors: genetics (multiple genes implicated—IRF6, MSX1, TBX22, FGFR1), maternal smoking, alcohol, diabetes, and folate deficiency. Diagnosis: antenatal ultrasound (around 20-week anomaly scan) or postnatal clinical examination (inspect palate with a tongue depressor). Complications: feeding difficulties (nasal regurgitation, poor weight gain), recurrent otitis media (Eustachian tube dysfunction—hearing loss), speech delay (velopharyngeal insufficiency—hypernasal speech, articulation errors), dental abnormalities (missing or malpositioned teeth), and maxillary hypoplasia. Management: multidisciplinary team (plastic surgeon, ENT, speech therapist, orthodontist, paediatrician, psychologist). Primary palatoplasty (repair of the palate) is performed at 9--12 months (before speech development); techniques: Furlow double-opposing Z-plasty, Bardach two-flap palatoplasty. Secondary surgery for velopharyngeal insufficiency (pharyngeal flap, sphincter pharyngoplasty). Alveolar bone grafting at 6--9 years. Orthognathic surgery for maxillary hypoplasia in adolescence. Speech therapy, audiology, and orthodontic treatment throughout childhood. Long-term outcomes: good speech and feeding outcomes, normal intelligence, and good quality of life with appropriate multidisciplinary care.

\medterm{Colectomy}-- Surgical resection of part or all of the colon. Right hemicolectomy: for cecal/ascending colon cancer. Left hemicolectomy: for descending/sigmoid cancer. Sigmoidectomy: for diverticulitis or sigmoid cancer. Total colectomy: for FAP, UC, or Lynch syndrome. Low anterior resection (LAR): for rectal cancer with anastomosis (coloanal anastomosis). Abdominoperineal resection (APR): for distal rectal cancer, permanent colostomy. Indications: colon cancer, diverticulitis, IBD (UC, Crohn), and ischemic colitis. Complications: anastomotic leak (most feared), hemorrhage, wound infection, and ileus.

\medterm{Colorectal Surgery}-- Colon cancer: right hemicolectomy (cecum to transverse), extended right (to mid-transverse), transverse colectomy, left hemicolectomy (splenic flexure to sigmoid), sigmoid colectomy, anterior resection (low anterior resection - LAR for upper/mid rectal), abdominoperineal resection (APR for low rectal with sphincter involvement), total mesorectal excision (TME) standard for rectal cancer. Lymphadenectomy: minimum 12 nodes. Rectal cancer: neoadjuvant chemoradiation (5-FU/capecitabine + 50.4 Gy) for T3/T4 or node-positive; total neoadjuvant therapy (TNT: induction chemotherapy followed by chemoradiation) for locally advanced; watch-and-wait for complete clinical response (cCR) after TNT. Anal cancer: squamous cell carcinoma; Nigro protocol (mitomycin C + 5-FU + 55-60 Gy radiation); surgery (APR) only for residual/recurrent disease. Diverticulitis: Hinchey classification (I: pericolonic abscess, antibiotics, percutaneous drainage if >4cm; II: pelvic abscess, percutaneous drainage; III: purulent peritonitis, Hartmann or resection with anastomosis + defunctioning stoma; IV: fecal peritonitis, Hartmann procedure). Elective sigmoid resection after 2 episodes of uncomplicated or 1 episode of complicated diverticulitis (current guidelines have shifted towards a more conservative, individualised approach, with surgery considered on a case-by-case basis for recurrent attacks, persistent complications, or immunosuppression).

\medterm{Craniotomy}-- Surgical opening of the skull (bone flap) for access to intracranial contents. Indications: brain tumor resection (glioblastoma, meningioma, mets), hematoma evacuation (SDH, EDH, ICH), aneurysm clipping, AVM resection, and epilepsy surgery (lobectomy). Stereotactic craniotomy: frameless neuronavigation (Stealth, BrainLab) for precise localization. Awake craniotomy: speech/language mapping for tumors in eloquent cortex. Complications: cerebral edema, hemorrhage (epidural, subdural, intraparenchymal), infection (osteomyelitis, meningitis), CSF leak, seizures, and neurologic deficits.

\medterm{Cryptorchidism}-- Failure of one or both testes to descend into the scrotum, the most common congenital genitourinary anomaly in boys (3\% at birth, 1\% at 1 year). 70\% of undescended testes are palpable (inguinal), 30\% non-palpable (intra-abdominal, absent, or atrophic). Spontaneous descent by 6 months occurs in about one-third. Complications: impaired fertility (bilateral: 75\% reduced, unilateral: 50\% reduced), testicular cancer risk (seminoma, 4-5x increased, even after orchiopexy), and psychological effects. Diagnosis: physical exam; ultrasound or laparoscopy for non-palpable. Treatment: orchiopexy (surgical placement in scrotum) between 6-18 months of age. Hormonal therapy (hCG) is less effective. Atrophic or absent testis: orchiectomy. Follow-up: annual exam and testicular self-exam.

\medterm{Debridement}-- The surgical removal of dead, damaged, or infected tissue from a wound to promote healing and prevent infection. Types: sharp/surgical debridement (scalpel, curette, scissors--most rapid and definitive), enzymatic debridement (collagenase, papain-urea), autolytic debridement (moist dressings allow the body's own enzymes to liquefy necrotic tissue), mechanical debridement (wet-to-dry dressings, irrigation, whirlpool), and biologic debridement (sterile maggots--Lucilia sericata larvae consume necrotic tissue). Selection depends on wound type, amount of necrotic tissue, infection status, and patient tolerance.

\medterm{Diabetic Foot Ulcer}-- A chronic wound of the foot in a patient with diabetes, resulting from peripheral neuropathy, peripheral artery disease, and impaired wound healing. Neuropathic ulcers: painless, at pressure points (plantar metatarsal heads, heel), punched-out appearance, surrounded by callus. Ischemic ulcers: painful, at toes or margins of foot, with pale, cool, hairless skin. Risk factors: poor glycemic control, long diabetes duration, prior amputation, neuropathy, PAD, and foot deformity. Diagnosis: monofilament testing, ankle-brachial index, vascular duplex ultrasound, and wound assessment. Treatment: offloading (total contact cast, CAM boot, wheelchair), debridement, infection control (antibiotics, surgical drainage for osteomyelitis), revascularization for ischemic ulcers, and advanced wound care (negative pressure, growth factors, bioengineered skin). Prevention: daily foot inspection, proper footwear, glycemic control, and podiatry referral.

\medterm{Dilation and Curettage (D\&C)}-- Surgical dilation of the cervix and curettage (scraping) of the uterine lining. Indications: incomplete or missed abortion (miscarriage management), therapeutic abortion, endometrial sampling (abnormal uterine bleeding, postmenopausal bleeding, endometrial hyperplasia/cancer), and removal of retained products of conception (RPOC). Complications: uterine perforation (1-3\%, higher in postmenopausal, retroverted uterus), cervical laceration, hemorrhage, infection, and Asherman syndrome (intrauterine adhesions from overzealous curettage). Now often replaced by hysteroscopy-guided procedures.

\medterm{Discectomy}-- Surgical removal of a herniated intervertebral disc fragment compressing a nerve root. Microscopic discectomy (most common): small incision, microscope, minimal muscle dissection. Indications: radicular pain (sciatica, cervical radiculopathy) refractory to 6-12 weeks of conservative therapy, progressive motor deficit, and cauda equina syndrome (surgical emergency). Success rate: 85-95\% for leg pain relief. Also performed via endoscopic, minimally invasive (Tubular/METRx, MIS-TLIF for associated instability), and open (rarely). Complications: dural tear (3-5\%), infection, hematoma, recurrent herniation, and nerve root injury.

\medterm{Drowning}-- Respiratory impairment from submersion/immersion in liquid. Drowning can be fatal or non-fatal. Pathophysiology: laryngospasm followed by aspiration of fluid, surfactant washout, pulmonary edema, and hypoxemia. Fresh water: washes out surfactant, causes alveolar collapse. Salt water: draws fluid into alveoli (pulmonary edema). Both cause severe hypoxemia, metabolic acidosis, and multiorgan failure. Risk factors: children 1-4 (bathtubs, pools), alcohol use, lack of supervision, and seizure disorders. Management: immediate rescue, CPR with rescue breaths (most important intervention), ICU care with ventilation, and treatment of aspiration pneumonia and ARDS. Prognosis: associated with submersion time, CPR delay, water temperature, and neurologic status on arrival.

\medterm{Electrical Injury}-- Tissue damage caused by passage of electrical current through the body. Severity depends on voltage, current (alternating vs. direct), pathway, duration, and resistance. High-voltage (>1000V): extensive internal tissue damage (muscle necrosis, compartment syndrome, rhabdomyolysis) with relatively small skin burns (entry/exit wounds). Cardiopulmonary effects: arrhythmias (ventricular fibrillation, asystole), cardiac arrest. Neurologic effects: spinal cord injury, peripheral nerve damage, cognitive deficits. Complications: acute kidney injury (myoglobinuria), compartment syndrome, fractures from falls, and delayed neurologic deficits. Diagnosis: ECG, CK, urine myoglobin, CT/MRI. Treatment: ABCs, cardiac monitoring, IV fluids with alkalinization (bicarbonate), surgical fasciotomy/debridement for compartment syndrome. Tetanus prophylaxis.

\medterm{Erectile Dysfunction}-- The persistent inability to achieve or maintain an erection sufficient for satisfactory sexual performance. Affects up to 50\% of men aged 40-70. Etiology: organic (vascular disease, diabetes, hypogonadism, neurological disorders, medications, Peyronie disease) and psychogenic (anxiety, depression, relationship issues). Vascular causes: impaired arterial inflow (atherosclerosis) or venous leak. Diagnosis: history, IIEF-5 questionnaire, nocturnal penile tumescence testing, duplex Doppler ultrasound, and hormonal profile. Treatment: PDE5 inhibitors (sildenafil, tadalafil, vardenafil) as first-line; second-line options: intracavernosal injections (alprostadil), vacuum erection devices, and penile prosthesis for refractory cases. Lifestyle modification (weight loss, exercise, smoking cessation) improves outcomes.

\medterm{Esophageal Surgery}-- GERD surgery: Nissen fundoplication (360° wrap, gold standard), partial wraps (Toupet 270°, Dor anterior 180-200°) for dysmotility. LINX device (magnetic sphincter augmentation) alternative. Achalasia: Heller myotomy (laparoscopic/robotic) + partial fundoplication (Dor); POEM (peroral endoscopic myotomy) endoscopic alternative; pneumatic dilation (recurrence 30-50\%). Esophageal cancer: adenocarcinoma (distal, Barrett's), squamous (proximal/mid, smoking/alcohol). Staging: EUS, PET-CT, laparoscopy for peritoneal disease. Neoadjuvant chemoradiation (CROSS regimen: carboplatin + paclitaxel + 41.4 Gy) for locally advanced (T2N+ or T3-4 any N). Surgery: Ivor Lewis (right thoracotomy + laparotomy, for mid/distal), McKeown (right thoracotomy + laparotomy + cervical, for proximal), transhiatal (blunt dissection, for distal, higher anastomotic leak rate). Reconstruction: gastric conduit (greater curve based on right gastroepiploic artery, most common), colon interposition (when stomach unavailable). Three-field lymphadenectomy (cervical, mediastinal, abdominal) for squamous cell carcinoma at or above the carina. Anastomotic techniques: stapled (EEA, 21-25 mm, end-to-side) or hand-sewn (two-layer running or interrupted, anastomotic leak rate 5-15\%). Wait time of 8-12 weeks after neoadjuvant therapy for surgical resection. Outcomes: 5-year survival for esophageal cancer is 15-25\% overall; 40-50\% for patients who achieve a pathological complete response (pCR) after neoadjuvant chemoradiation. Salvage oesophagectomy for persistent or recurrent disease after definitive chemoradiation (50.4 Gy + chemotherapy) has higher morbidity (30-50\% complication rate, especially anastomotic leak and pulmonary complications).

\medterm{Excision}-- The surgical removal of a piece of tissue, an organ, or a lesion. Excisional biopsy removes the entire lesion with a margin of normal tissue, used for definitive diagnosis and treatment of small tumors (suspected melanoma, breast lump, lymph node). Wide local excision removes the lesion with an additional margin of healthy tissue for oncologic clearance. In biochemistry, excision repair: the removal of damaged DNA segments. Surgical principles: hemostasis, closure method (primary, secondary, or tertiary intention), and specimen orientation for pathology. Margin status (clear, close, or involved) is the most important prognostic factor for local recurrence in cancer surgery.

\medterm{Fasciotomy}-- Surgical incision through the fascia to relieve compartment syndrome. Indications: acute compartment syndrome (trauma, reperfusion, burns, crush injury). Classic signs: pain out of proportion, pain with passive stretch, paresthesias, pallor, pulselessness (late). Diagnosis: clinical, compartment pressure >30 mmHg OR delta pressure (diastolic BP minus compartment pressure) <30 mmHg. Leg: 4 compartments (anterior, lateral, superficial posterior, deep posterior) via two incisions. Forearm: volar (Henry approach) and dorsal. Foot, thigh, hand, abdomen (for abdominal compartment syndrome) also possible. Fasciotomy wounds are left open for 48-72 hours, then delayed primary closure or skin grafting.

\medterm{Femoral Hernia}-- A hernia of abdominal contents through the femoral ring, into the femoral canal, medial to the femoral vessels and below the inguinal ligament. More common in women (3:1) and more prone to strangulation (40\%) than inguinal hernias due to narrow, rigid femoral ring. Presents as a small, non-reducible mass below and lateral to the pubic tubercle (distinguishing from inguinal hernia). Diagnosis: clinical; ultrasound/CT for equivocal cases. Treatment: surgical repair (McVay, mesh plug, or laparoscopic TEP/TAPP) due to high strangulation risk. Emergency surgery for strangulated or incarcerated femoral hernia.

\medterm{Frostbite}-- Tissue injury caused by freezing of skin and underlying tissues. Severity grading: first degree (superficial, numbness, erythema), second degree (blisters, clear fluid, pain), third degree (hemorrhagic blisters, necrosis, full-thickness skin loss), fourth degree (deep tissue necrosis, mummification, autoamputation). Common sites: fingers, toes, nose, ears, cheeks. Predisposing factors: cold exposure, wind chill, vasoconstriction, immobility, alcohol, smoking, diabetes. Pathophysiology: ice crystal formation, endothelial injury, thrombosis, and reperfusion injury upon rewarming. Diagnosis: clinical. Treatment: rapid rewarming in water 37-39°C (100-104°F) for 15-30 minutes; ibuprofen, tetanus prophylaxis, wound care, and delayed amputation (march of demarcation). Thrombolytic therapy (tPA) within 24 hours for severe cases to reduce amputation rate.

\medterm{Gallbladder Cancer}-- Aggressive, often diagnosed incidentally after cholecystectomy for cholelithiasis (1-2\% of cholecystectomies). Risk factors: gallstones (chronic inflammation), porcelain gallbladder (calcified wall, 10-25\% cancer risk), anomalous pancreaticobiliary junction, gallbladder polyps >1cm, typhoid carrier state. T staging: T1a (lamina propria), T1b (muscularis), T2 (perimuscular connective tissue), T3 (perforates serosa or invades liver/adjacent organ), T4 (main portal vein/hepatic artery, extensive). Lymphatic spread: cystic, pericholedochal, hilar, peripancreatic, periduodenal, celiac, retroperitoneal nodes. Surgical management: T1a: simple cholecystectomy alone is sufficient (no lymphadenectomy required, 5-year survival >95\%); T1b: extended cholecystectomy (cholecystectomy + en bloc wedge resection of liver segments IVb and V + portal lymphadenectomy) is the minimum recommended operation (5-year survival 80-85\%); T2: extended cholecystectomy + portal lymphadenectomy (5-year survival 50-60\%); T3: extended cholecystectomy + portal lymphadenectomy ± bile duct resection/right hepatectomy (5-year survival 20-30\%); T4: unresectable (palliative: biliary stenting, chemotherapy). Incidental gallbladder cancer discovered after simple cholecystectomy: refer to a hepatobiliary center for completion radical cholecystectomy (wedge resection of gallbladder bed, portal lymphadenectomy, ± bile duct excision and reconstruction with Roux-en-Y hepaticojejunostomy) within 4 weeks of the initial cholecystectomy. Chemotherapy: gemcitabine + cisplatin for advanced disease (ABC-02 and BILCAP trials). Prognosis: overall 5-year survival 5-15\%, worse if lymph node involvement (5-year survival <10\% with nodal metastases). Adjuvant chemotherapy with capecitabine (BILCAP) or gemcitabine-based regimens may improve survival after curative resection.

\medterm{Gastric Surgery}-- Gastric cancer: D2 lymphadenectomy standard (stations 1-12a, 12b). Gastrectomy types: total, distal subtotal (Billroth I gastroduodenostomy, Billroth II gastrojejunostomy, Roux-en-Y), proximal, pylorus-preserving. Reconstruction: Roux-en-Y preferred (prevents bile reflux). Early gastric cancer (T1a): ESD (endoscopic submucosal dissection) if no LVI, well-differentiated, <2cm (expanded criteria). Neoadjuvant: FLOT (5-FU, leucovorin, oxaliplatin, docetaxel) superior to surgery alone. Palliative: systemic chemotherapy, palliative resection for bleeding/obstruction/perforation, stenting for GOO (gastric outlet obstruction). GIST: imatinib (neoadjuvant for downstaging, adjuvant if high risk: >10cm, mitotic >10/HPF, rupture), surgery (R0, avoid rupture). MALT lymphoma: H. pylori eradication first; only treat with surgery or chemoradiation if not responsive. Perforated peptic ulcer: omental patch repair (Graham patch: interrupted 2-0 silk sutures through the full thickness of each side of the perforation, tied over a pedicle of omentum, which is then secured in place with a second set of sutures), laparoscopic (preferred if hemodynamically stable, no contamination, surgeon experience), open if peritonitis. Intraoperative decision-making: if the perforation is >2cm, there is a suspicion of malignancy, or the edges are oedematous and friable, a more extensive resection may be required (distal gastrectomy with Billroth II or Roux-en-Y reconstruction). Post-operative: PPI infusion (omeprazole 80 mg bolus + 8 mg/hour), H. pylori testing and eradication (triple therapy: PPI + amoxicillin + clarithromycin or metronidazole, 14-day course), return to oral intake after resolution of ileus.

\medterm{Grafting}-- The surgical transfer of tissue from one site to another, or from one individual to another. Types: (1) Autograft: same individual (gold standard, e.g., skin graft for burns, bone graft for non-union, saphenous vein graft for CABG). (2) Allograft: same species, different individual (cadaveric bone, cornea, heart valve, kidney). (3) Xenograft: different species (porcine heart valve, porcine/bovine skin for burns). (4) Isograft: identical twin (syngeneic). (5) Synthetic graft: prosthetic material (Dacron vascular graft, Gore-Tex, metal implants, polyethylene for joint replacement). Skin grafts: split-thickness (epidermis + part of dermis) or full-thickness (epidermis + full dermis). Flaps: tissue transferred with its blood supply intact (pedicled or free). Complications: graft rejection (allografts, xenografts), infection, haematoma, seroma, graft failure (ischemia), contracture. Immunosuppression (tacrolimus, mycophenolate, corticosteroids) is required for allograft organ transplantation.

\medterm{Heimlich Manoeuvre}-- An emergency procedure for relieving upper airway obstruction caused by a foreign body (choking). The maneuver is performed on a conscious, choking adult who cannot speak, cough effectively, or breathe. Technique: the rescuer stands behind the victim, wraps their arms around the victim's waist, makes a fist with one hand (thumb side against the victims abdomen, midline above the navel and well below the xiphoid process), grasps the fist with the other hand, and performs upward, inward thrusts (abdominal thrusts). For obese or pregnant victims, chest thrusts (similar position but on the lower sternum) are used instead. The maneuver abruptly increases intra-abdominal and intrathoracic pressure, forcing air out of the lungs and creating a cough-like force to expel the obstructing object. For a choking infant (under 1 year): hold the infant face-down along the rescuer's forearm, with the head lower than the trunk, and deliver 5 back blows between the shoulder blades, followed by 5 chest thrusts (face-up, two fingers on lower sternum). The Heimlich maneuver is contraindicated for: infants under 1 year (use back blows and chest thrusts), unconscious victims (start CPR), and cases where the obstruction is not complete (partial obstruction with good air exchange, encourage coughing). After successful expulsion, the patient should have a medical evaluation to exclude complications (esophageal rupture, gastric rupture, pneumothorax, abdominal organ injury).

\medterm{Hemicolectomy}-- See Colectomy.

\medterm{Hemorrhoidectomy}-- Surgical excision of symptomatic hemorrhoids. Indications: Grade III-IV internal hemorrhoids (prolapsed, irreducible), external hemorrhoids with uncontrolled bleeding or recurrent thrombosis, mixed hemorrhoids. Techniques: closed (Ferguson, primary closure), open (Milligan-Morgan, leave wounds open), and stapled hemorrhoidopexy (PPH, for circumferential prolapse, lower pain but higher recurrence). Laser, bipolar, and harmonic scalpel techniques also available. Postoperative pain is significant (particularly with open technique). Recovery: 2-4 weeks. Complications: pain, bleeding, urinary retention, anal stenosis, and incontinence.

\medterm{Hepatic Surgery}-- Liver resection: Couinaud segments. Indications: colorectal liver metastases (CRLM, up to 70\% 5-year survival with R0), HCC (Child-Pugh A, portal pressure <10 mmHg, future liver remnant >20\% normal, >40\% cirrhotic), cholangiocarcinoma, benign tumors (adenoma >5cm, symptomatic hemangioma, FNH diagnostic uncertainty). Preoperative: PVE (portal vein embolization) for FLR <20\%/40\%; ALPPS (associating liver partition and portal vein ligation) for rapid hypertrophy (2 weeks). Techniques: open, laparoscopic, robotic. Pringle manoeuvre (clamp hepatoduodenal ligament, 15-20 min on, 5 min off, ischemic preconditioning may reduce ischemia-reperfusion injury). Low central venous pressure (CVP <5 mmHg) anesthesia. Parenchymal transection: crush-clamp (Kelly clamp), CUSA (cavitron ultrasonic surgical aspirator), LigaSure, Harmonic scalpel. Hepatic vein management: extrahepatic control (right, middle, left suprahepatic veins) or intrahepatic approach (anterior approach with hanging manoeuvre). Complications: bile leak (3-10\%), PHLF (post-hepatectomy liver failure, defined by ISGLS peak bilirubin >50 mumol/L + INR >1.5 on post-operative day 5; management: medical support, albumin, diuretics, TIPS, or liver transplantation in severe cases), hemorrhage, abscess, ascites. Non-resectional techniques: RFA (radiofrequency ablation), MWA (microwave ablation), IRE (irreversible electroporation), SBRT (stereotactic body radiotherapy), TACE (transarterial chemoembolization), and TARE (transarterial radioembolization with yttrium-90), each indicated for specific tumor types and locations where surgical resection is not feasible. Liver transplantation: Milan criteria for HCC (single \(\le\)5 cm or up to 3 nodules each \(\le\)3 cm, no extrahepatic disease, no major vascular invasion) with 5-year survival >70\%.

\medterm{Hernia}-- The protrusion of an organ or tissue through an abnormal opening in the wall of the cavity that normally contains it. Most hernias occur in the abdominal wall. Types: (1) Inguinal hernia (most common): indirect (through the deep inguinal ring, lateral to inferior epigastric vessels, congenital or acquired, more common in younger men) vs direct (through Hesselbach triangle, medial to inferior epigastric vessels, acquired, older men). (2) Femoral hernia: through the femoral ring below the inguinal ligament, more common in women, higher risk of strangulation (40\%) vs inguinal (10\%). (3) Umbilical hernia: through the umbilical ring, common in infants (usually resolves spontaneously by age 2-4) and adults (obesity, pregnancy, ascites). (4) Incisional (ventral) hernia: through a surgical scar, risk factors: wound infection, obesity, smoking, COPD. (5) Hiatus hernia: see separate entry. (6) Epigastric hernia: through the linea alba. (7) Spigelian hernia: through the linea semilunaris. (8) Obturator hernia: through the obturator foramen (rare, elderly thin women). (9) Lumbar hernia: through the lumbar triangle. (10) Diaphragmatic hernia: through the diaphragm. Diagnosis: history (lump or bulge, worse on standing/straining, cough impulse), examination (inspection: visible bulge; palpation: cough impulse; reducibility). Complications: irreducibility (cannot be reduced), obstruction (bowel obstruction), strangulation (vascular compromise--ischemia, gangrene, surgical emergency). Treatment: surgical repair (open or laparoscopic; tension-free mesh repair is standard--Lichtenstein for inguinal, laparoscopic transabdominal preperitoneal TAPP or totally extraperitoneal TEP). Watchful waiting is an option for minimally symptomatic inguinal hernias in men.

\medterm{Hernia Repair (Herniorrhaphy)}-- Surgical repair of a hernia (inguinal, femoral, umbilical, incisional, hiatal). Inguinal hernia repair: open (Lichtenstein tension-free mesh, Shouldice, Bassini) or laparoscopic (TEP: totally extraperitoneal, TAPP: transabdominal preperitoneal). Mesh reinforcement reduces recurrence (1-3\% for mesh vs. 5-15\% for primary tissue repair). Indications: symptomatic hernia, incarceration/strangulation risk, and groin pain. Emergency repair for strangulated hernia (bowel resection may be needed). Complications: recurrence, chronic groin pain (5-10\%, especially open), mesh infection, seroma, hematoma, and cord injury (vas deferens, testicular vessels).

\medterm{Hiatus Hernia}-- The protrusion of part of the stomach through the esophageal hiatus of the diaphragm into the chest cavity. Types: (1) Sliding hernia (type I, 95\%): the gastro-esophageal junction (GOJ) and a portion of the gastric cardia slide above the diaphragm through the hiatus. The GOJ is displaced cephalad, and the angle of His is lost, predisposing to GORD. Associated with: obesity, pregnancy, increased intra-abdominal pressure. (2) Para-esophageal hernia (type II): the GOJ remains in its normal position below the diaphragm, but a portion of the gastric fundus herniates alongside the esophagus. (3) Mixed (type III): both GOJ and fundus herniate. (4) Type IV: herniation of other abdominal organs (colon, omentum, spleen) through the hiatus. Para-esophageal hernias (types II-IV) can incarcerate or strangulate. Symptoms: many sliding hernias are asymptomatic; associated with GORD (heartburn, regurgitation, dysphagia, chest pain, chronic cough). Large para-esophageal hernias may cause: postprandial fullness, chest pain, dyspnoea, dysphagia, vomiting, and rarely acute gastric volvulus (Borchardt triad: severe epigastric pain, retching without vomiting, inability to pass a nasogastric tube--surgical emergency). Diagnosis: barium swallow (GOJ above diaphragm), upper endoscopy (GOJ displacement, Los Angeles classification of oesophagitis), high-resolution manometry. Management: asymptomatic sliding hernias require no treatment; symptomatic GORD: PPI therapy, lifestyle measures; para-esophageal hernias with symptoms or risk of volvulus: laparoscopic Nissen fundoplication with hernia reduction and crural repair (hiatoplasty).

\medterm{Hypospadias}-- A congenital anomaly in which the urethral meatus opens on the ventral (underside) of the penis, proximal to the glans. Incidence: 1 in 250 male births. Classification: distal (glanular, coronal — 70\%), midshaft (10\%), proximal (penoscrotal, scrotal, perineal — 20\%). Associated with chordee (ventral penile curvature), undescended testis, and inguinal hernia. Most cases are idiopathic; genetic and endocrine factors are implicated. Diagnosis: clinical exam — note the position of the meatus and presence of chordee. Do not circumcise before surgical repair (foreskin may be needed for urethroplasty). Surgical repair (urethroplasty) at 6-18 months: single-stage (TIP / Snodgrass repair for distal) or two-stage (for proximal). Complications: fistula, stricture, and residual chordee.

\medterm{Hysterectomy}-- Surgical removal of the uterus. Total hysterectomy: removal of uterus and cervix. Subtotal (supracervical): removes uterus, leaves cervix. Radical hysterectomy: uterus, cervix, parametrium, upper vagina, and pelvic lymph nodes (for cervical cancer). Routes: abdominal (open, most common for large uteri/malignancy), vaginal (for prolapse, small uterus), laparoscopic (TLH, less pain, shorter stay), and robotic (da Vinci, improved dexterity). Indications: uterine fibroids (most common), endometriosis, adenomyosis, uterine prolapse, abnormal uterine bleeding, endometrial cancer, cervical cancer, and ovarian cancer. Complications: bleeding, infection, ureter/bladder/bowel injury, and vaginal cuff dehiscence.

\medterm{Incisional Hernia}-- A hernia through a previous surgical incision scar, the most common long-term complication of abdominal surgery (10-20\% of laparotomies). Risk factors: wound infection, obesity, smoking, malnutrition, COPD, and multiple previous surgeries. Present as a bulge at or near the incision site, worse with standing or straining. Can cause pain, obstruction, and strangulation. Diagnosis: clinical; CT for evaluation of size, number of defects, and contents. Treatment: surgical repair with mesh (open sublay, onlay, or laparoscopic IPOM) to reduce recurrence. Component separation for large/complex hernias.

\medterm{Inguinal Hernia}-- Protrusion of abdominal contents through the inguinal canal. Indirect (lateral to inferior epigastric vessels, through deep inguinal ring): congenital, patent processus vaginalis; most common overall (60-70\%), male predominance. Direct (medial to inferior epigastric vessels, through Hesselbach triangle): acquired, weakness of transversalis fascia; older men. Pantaloon: both indirect and direct. Femoral hernia (below inguinal ligament, through femoral canal): more common in women, right-sided predominance. Clinical: groin bulge (medial to pubic tubercle, below inguinal ligament, irreducible, may cause obstruction or strangulation). Imaging: ultrasound (sensitivity 95\%, Valsalva enhancement increases sensitivity) or MRI if diagnosis uncertain. Management: elective repair (symptomatic; open or laparoscopic: TEP - totally extraperitoneal, TAPP - transabdominal preperitoneal). Mesh (polypropylene, lightweight macroporous) preferred over tissue repair (Bassini, Shouldice, McVay, Desarda, Lichtenstein) due to <5\% recurrence. Lichtenstein: open tension-free mesh onlay repair, gold standard for open inguinal hernia repair. Laparoscopic: lower recurrence for bilateral and recurrent hernias, faster recovery, less chronic pain (in experienced hands, 1-2\% for TEP). Emergency surgery for obstruction/strangulation: McBurney or Lotheissen approach (low midline incision for control of incarcerated bowel), viable bowel → herniorrhaphy, nonviable bowel → resection + herniorrhaphy. Femoral hernia repair: McEvedy approach (retroperitoneal, high approach, through a lower midline or transverse lower abdominal incision) or laparoscopic TEP/TAPP; locking sutures or mesh plug (Bard plug, PerFix plug) or open suture repair (Moschcowitz type, a suture repair of the femoral canal). Complications: recurrence (1-10\%, higher with tissue repairs and laparoscopic learning curve), chronic groin pain (10-20\%, higher with open, meshoma, nerve entrapment, triple neurectomy for refractory cases), seroma, haematoma, infection (mesh infection <1\%, may require mesh removal), testicular atrophy or ischemic orchitis (1-2\% during open repair for recurrent hernias or large scrotal hernias), and urinary retention (5-10\%, transient) in older male patients after general or spinal anesthesia.

\medterm{Jejunostomy}-- See Feeding Jejunostomy.

\medterm{Laceration}-- A wound caused by tearing of body tissue, typically caused by trauma from a sharp object (knife, glass, metal) or blunt force that splits the skin. Lacerations involve the skin and may extend into subcutaneous tissue, muscle, tendon, nerve, or blood vessels. Assessment: location, length, depth, contamination, foreign bodies, neurovascular status (sensation, motor function, capillary refill, distal pulses), and involvement of deeper structures (tendon function, bone exposure). Management: (1) hemostasis: direct pressure, elevation. (2) Wound cleaning: irrigation with normal saline (high-pressure, 250-500 mL, reduces bacterial load). (3) anesthesia: local infiltration with lidocaine 1-2\% (with or without adrenaline, max 4.5 mg/kg plain, 7 mg/kg with adrenaline) or topical lidocaine-prilocaine cream for wound margins. (4) Exploration: assess for foreign bodies, tendon/nerve/vascular injury. (5) Debridement: excision of devitalised tissue, ragged wound edges (freshening). (6) Closure: primary closure (sutures--simple interrupted, continuous/mattress, subcuticular; tissue adhesive; adhesive strips; staples) for clean, non-infected wounds; delayed primary closure (3-5 days) for contaminated wounds; healing by secondary intention for infected or high-risk wounds. Tetanus prophylaxis: clean wound--tetanus toxoid if last dose >10 years; contaminated wound--tetanus toxoid + immunoglobulin if last dose >5 years or unknown. Antibiotics: not routinely indicated for simple lacerations; consider for contaminated wounds, human/animal bites, immunocompromised patients. Sutures removal timing: face 3-5 days, scalp 7 days, trunk/arms 7-10 days, legs 10-14 days. Complications: wound infection (2-5\%), dehiscence, hypertrophic scarring/keloid, retained foreign body, nerve damage.

\medterm{Laminectomy}-- Surgical removal of the lamina (posterior vertebral arch) to decompress the spinal canal. Indications: spinal stenosis (lumbar and cervical), herniated disc (with central stenosis), and spinal tumor. Decompression of the entire spinal canal. Often combined with discectomy (lumbar) or foraminotomy (nerve root decompression). May require fusion if instability is present (laminectomy with fusion, PLF). Complications: CSF leak (dural tear), spinal instability (post-laminectomy kyphosis), epidural hematoma, infection, and neurologic injury.

\medterm{Large Bowel Obstruction}-- Less common than SBO. Causes: colorectal cancer (60\%), diverticular stricture (20\%), volvulus (sigmoid 60\%, cecal 40\%), fecal impaction, Ogilvie syndrome (acute colonic pseudo-obstruction), ischemic stricture, endometriosis, radiation stricture. Clinical: abdominal distension (prominent), constipation/obstipation, less vomiting than SBO, pain may be less severe. Imaging: abdominal X-ray (dilated colon >6cm, cecum >9cm predicts perforation), CT with contrast (transition point, cause, closed loop obstruction, pneumatosis intestinalis indicating ischemia, portal venous gas, pneumoperitoneum - all signs of ischemia or perforation). Water-soluble contrast enema can differentiate pseudo-obstruction (Ogilvie) from mechanical obstruction (transition point, cut-off sign). Management: NPO, NGT, IV fluids, correct electrolyte abnormalities. Specific management depends on the cause. Sigmoid volvulus: endoscopic detorsion (rigid sigmoidoscopy with placement of a flatus tube, success rate 80-90\%) + elective sigmoid colectomy (resection with primary anastomosis) during the same admission (recurrence 40-60\% without surgery). Cecal volvulus: colonoscopic detorsion (less successful, 30-40\%) + right hemicolectomy. Colorectal cancer: segmental colectomy with primary anastomosis (elective); emergency: resection with stoma (Hartmann) or resection with on-table colonic lavage and primary anastomosis (controversial, increased leak risk). Stenting (self-expanding metallic stent, SEMS) as bridge to surgery (if no perforation, no peritonitis, and the tumor is accessible, allows for medical optimisation and elective resection with primary anastomosis avoiding a stoma in 80-90\% of cases) or as definitive palliation (for metastatic disease, high surgical risk). Ogilvie syndrome (acute colonic pseudo-obstruction): conservative (NPO, NGT, correction of electrolytes, mobilisation), neostigmine (2 mg IV, continuous cardiac monitoring for bradycardia and bronchospasm, only in the absence of mechanical obstruction confirmed by CT or water-soluble contrast enema; response within 30 minutes in 60-90\% of patients), colonoscopic decompression (suction of air, aspiration of liquid stool) for neostigmine failure, cecostomy (percutaneous or surgical) for recurrence. Contraindications to neostigmine: mechanical obstruction, bradycardia, hypotension, bronchial asthma, peptic ulcer disease, pregnancy. Surgical intervention (laparotomy with colectomy, subtotal colectomy with terminal ileostomy and Hartmann closure of the distal stump) for perforation, ischemia, or failure of all conservative measures.

\medterm{Lobectomy}-- Surgical removal of a lobe of an organ. Most commonly lung lobectomy (removal of one lobe of the lung, performed for lung cancer, bronchiectasis, abscess, or benign tumor) or temporal lobectomy (removal of the anterior temporal lobe for drug-resistant temporal lobe epilepsy). Lung lobectomy: the standard surgical treatment for early-stage (stage I-II) non-small cell lung cancer (NSCLC). Approaches: open thoracotomy (posterolateral) or video-assisted thoracoscopic surgery (VATS, minimally invasive, less post-operative pain, shorter hospital stay, equivalent oncologic outcomes). The remaining lung lobes expand to fill the thoracic cavity. Pulmonary function loss corresponds to the number of segments removed (each lobe contains 2-5 bronchopulmonary segments). Compensatory hyperinflation of remaining lung occurs over weeks. Temporal lobectomy: anterior temporal lobectomy (ATL) for epilepsy involves resection of 3-4 cm of the anterior temporal neocortex plus the amygdala and anterior 2-3 cm of the hippocampus. Seizure-free outcome: 60-70\% in well-selected patients with mesial temporal sclerosis. Complications: visual field deficit (superior quadrantanopia, due to optic radiation Meyer loop injury, 50-80\%), language deficits (dominant hemisphere), memory impairment. Other lobectomies: hepatic lobectomy (liver resection, for cancer or living donor), thyroid lobectomy (removal of one thyroid lobe, for thyroid nodule/cancer), and cerebral lobectomy (removal of a whole brain lobe, for severe trauma or tumor).

\medterm{Mastectomy}-- Surgical removal of the breast. Total (simple) mastectomy: removal of breast tissue, areola, and nipple. Modified radical mastectomy (MRM): total mastectomy + axillary lymph node dissection (levels I-III), for invasive breast cancer. Skin-sparing mastectomy: preserves breast skin envelope for immediate reconstruction (implant or autologous). Nipple-sparing mastectomy: preserves nipple-areolar complex (for small, peripherally located tumors). Indications: breast cancer (locally advanced, multicentric, inflammatory), and risk-reducing in BRCA carriers (prophylactic bilateral mastectomy). Sentinel lymph node biopsy (SLNB) replaced ALND for node-negative patients. Reconstruction: implant-based or autologous (TRAM, DIEP, latissimus dorsi flap).

\medterm{Necrotizing Fasciitis}-- A life-threatening soft tissue infection characterized by rapid necrosis of the subcutaneous fascia and subcutaneous tissue, with relative sparing of the overlying skin in early stages. It is a surgical emergency requiring immediate wide surgical debridement and broad-spectrum antibiotics. Type I (polymicrobial, 60-75 percent of cases) occurs in diabetic, immunocompromised, or post-surgical patients and involves aerobic and anaerobic bacteria (E. coli, Klebsiella, Bacteroides, Peptostreptococcus, Clostridium species). Type II (monomicrobial, 15-30 percent) is caused by Group A Streptococcus (Streptococcus pyogenes) with or without Staphylococcus aureus, often in healthy individuals, and may present with toxic shock syndrome. Type III (gram-negative, marine-related: Vibrio vulnificus, Aeromonas hydrophila) occurs after exposure to seawater or freshwater, especially in patients with chronic liver disease. Type IV (fungal: Candida, Mucor, Aspergillus) occurs in severely immunocompromised patients and neutropenic patients. LRINEC score (Laboratory Risk Indicator for Necrotizing Fasciitis) stratifies risk: CRP >150 mg/L + WBC >15,000 + Hb <11 g/dL + Na <135 mmol/L + Cr >1.6 mg/dL + glucose >10 mmol/L, with a score \(\ge\)6 raising suspicion (positive predictive value 92\%) and \(\ge\)8 strongly predictive (positive predictive value 93-99\%). Management: aggressive IV fluid resuscitation, broad-spectrum antibiotics (piperacillin-tazobactam 4.5 g IV q6h or meropenem 1 g IV q8h + clindamycin 600-900 mg IV q8h to suppress toxin production, particularly for streptococcal toxic shock syndrome), urgent surgical debridement (fasciotomy with excision of all devitalised tissue, to the level of healthy bleeding muscle, re-exploration within 24-48 hours is mandatory). IVIG is used as adjunctive therapy for streptococcal toxic shock syndrome (2 g/kg loading dose, then 0.4 g/kg/day for 3-5 days). Hyperbaric oxygen (HBO) therapy is controversial but may be beneficial as an adjunct. Post-operative care includes repeat debridement such that the wound is thoroughly debrided within 24-72 hours, negative pressure wound therapy (NPWT) to promote granulation, and delayed primary closure or split-thickness skin grafting (STSG) once the wound is clean. Mortality: 20-40\%, >50\% if septic shock at presentation. Amputation rate is 15-25\%. Poor prognostic factors include advanced age (>60), diabetes mellitus, delay in surgery (>24 hours from admission to first debridement), involvement of the trunk, and streptococcal toxic shock syndrome. Renal failure, immune compromise, LRINEC score \(\ge\)8 and the presence of bullae or skin necrosis are also poor prognostic signs.

\medterm{Nephrectomy}-- Surgical removal of a kidney. Radical nephrectomy: kidney + Gerota fascia + adrenal gland (for renal cell carcinoma). Partial nephrectomy (nephron-sparing): for small renal tumors (<7 cm, renal function preservation). Simple nephrectomy: for benign disease (non-functioning kidney, chronic infection, massive hydronephrosis). Laparoscopic/robotic approach is standard (less pain, shorter stay). Indications: renal cell carcinoma (most common), trauma, and living donor nephrectomy. Complications: bleeding, infection, pneumothorax, splenic or pancreatic injury (left), and vascular injury.

\medterm{Nephrostomy Tube}-- A percutaneously placed catheter in the renal pelvis to drain urine, bypassing a ureteral obstruction. Indications: ureteral obstruction (calculus, tumor, stricture) causing hydronephrosis, pyonephrosis (infected, obstructed kidney — emergency drainage), and as a temporary measure before definitive surgery. Placed by interventional radiology under ultrasound/CT guidance. Sizes: 8-14 Fr. Output: normal 1-2 L/day. Tube is secured to the skin and connected to a drainage bag. Changed q6-12 weeks. Complications: bleeding (hematuria, perirenal hematoma), infection, tube dislodgement/obstruction, and adjacent organ injury (colon, pleura).

\medterm{Open-Heart Surgery}-- Any surgical procedure involving opening the chest wall (sternotomy or thoracotomy) and operating on the heart, usually with cardiopulmonary bypass (the heart--lung machine). Procedures: coronary artery bypass grafting (CABG—most common), valve repair/replacement (aortic, mitral), septal defect repair (ASD, VSD), heart transplantation, and ventricular assist device implantation. Cardiopulmonary bypass drains venous blood, oxygenates it, and returns it to the arterial system, allowing the surgeon to operate on a still, bloodless heart. Cardioplegia solution (high potassium) stops the heart in diastole, reducing myocardial oxygen demand during ischemia. Risks: bleeding, infection (mediastinitis, sternal wound), stroke (aortic manipulation—embolic), acute kidney injury, post-pericardiotomy syndrome, and cognitive dysfunction ('pump head'—typically transient). Minimally invasive approaches (mini-thoracotomy, robotic-assisted) reduce recovery time and blood loss. Off-pump CABG avoids cardiopulmonary bypass but has similar outcomes in selected patients.

\medterm{Organ Transplants}-- The surgical transfer of an organ from a donor to a recipient with end-stage organ failure. Solid organ transplants: kidney (most common, >90,000 per year worldwide), liver, heart, lung, pancreas, and intestine. Donor types: deceased donor (brainstem death or donation after circulatory death—DCD), living donor (kidney, partial liver, partial lung). Immunosuppression: induction (anti-thymocyte globulin, basiliximab) and maintenance (calcineurin inhibitors—tacrolimus, cyclosporine; antiproliferatives—mycophenolate mofetil; corticosteroids). Rejection: hyperacute (minutes—preformed antibodies, ABO incompatibility), acute (days to months—T-cell mediated, treatable with high-dose steroids), and chronic (months to years—multifactorial, progressive graft dysfunction). Organ allocation is managed by national bodies (NHSBT in the UK, UNOS in the US). Waiting list mortality is driven by organ shortage. Xenotransplantation (pig organs—genetically modified to reduce hyperacute rejection) is emerging as a potential solution.

\medterm{Pancreatic Neuroendocrine Tumors}-- Pancreatic neuroendocrine tumors (pNETs): functional (insulinoma, gastrinoma, glucagonoma, VIPoma, somatostatinoma) or non-functional (60-70\%). WHO 2017 grading: NET G1 (Ki-67 <3\%, mitotic <2/10 HPF), NET G2 (Ki-67 3-20\%, mitotic 2-20/10 HPF), NEC G3 (Ki-67 >20\%, mitotic >20/10 HPF). Functional: insulinoma (hypoglycemia, Whipple triad, 90\% benign, enucleation/distal pancreatectomy), gastrinoma (Zollinger-Ellison syndrome, ulcers, diarrhea, MEN1 association, duodenal > pancreatic, surgery if localized and resectable, PPI (gastrinoma: high-dose PPI, 80 mg omeprazole daily, can be increased to 120 mg daily for symptom control), octreotide for symptom control), glucagonoma (necrolytic migratory erythema, diabetes, weight loss, DVT, distal pancreatectomy + octreotide), VIPoma (WDHA syndrome: watery diarrhea, hypokalaemia, achlorhydria, VIP levels >75 pmol/L, octreotide, surgery), somatostatinoma (diabetes, gallstones, steatorrhea, surgery). Non-functional: 60-70\%, large >2cm at diagnosis, no hormonal syndrome, surgery (resection or enucleation if <2cm and favourable location). pNETs in MEN1: multiple, less aggressive, surgery for >2cm or functioning tumors in the pancreatic head; enucleation preferred for superficial tumors while deeper lesions may require pancreaticoduodenectomy; observation for small (<2cm) non-functioning pNETs with surveillance every 6-12 months (in MEN1 and sporadic cases). Syndromic associations: MEN1 (multiple pNETs, parathyroid adenomas, pituitary adenomas, carcinoid tumors), VHL (serous cystadenomas, pNETs, renal cell carcinoma, phaeochromocytoma, retinal angiomas, cerebellar haemangioblastomas), tuberous sclerosis, NF1. Prognosis: 5-year survival 90\% for well-differentiated localised NET, 30-50\% for metastatic (grade-dependent; Ki-67 is the most important independent prognostic factor). WHO classification 2019: well-differentiated neuroendocrine tumor (NET G1, G2, G3) and poorly differentiated neuroendocrine carcinoma (NEC G3, small cell or large cell type), with significant difference in prognosis and management (NET G3 is treated as an indolent well-differentiated NET rather than aggressive NEC G3).

\medterm{Pancreatic Surgery}-- Pancreaticoduodenectomy (Whipple): pancreatic head, duodenum, gallbladder, distal stomach (pylorus-preserving standard), common bile duct; reconstruction: pancreaticojejunostomy (duct-to-mucosa or invagination), hepaticojejunostomy, gastrojejunostomy/duodenojejunostomy. Distal pancreatectomy: body/tail; spleen-preserving (Kimura, Warshaw) or splenectomy. Central pancreatectomy: neck tumors, preserve head and tail. Enucleation: small pNETs/insulinomas away from duct. Complications: POPF (pancreatic fistula; ISGPS definition: amylase >3x upper limit of normal in drain fluid on or after post-operative day 3; biochemical leak (grade A, no change in management) vs grade B (requires change in management, e.g., prolonged drainage, total parenteral nutrition, somatostatin analogues, antibiotics, percutaneous drainage) vs grade C (major change in management, ICU stay, reoperation, or death); management: NPO, TPN or enteral nutrition via nasojejunal tube if possible, octreotide (somatostatin 250-500 mcg subcutaneously TDS after gastrectomy or after Whipple procedure to reduce exocrine pancreatic secretion, though evidence is mixed and routine use is not recommended), ERCP with pancreatic duct stenting for persistent high-output fistula; percutaneous drainage of fluid collections; reoperation for uncontrolled fistula or bleeding), DGE (delayed gastric emptying, ISGPS definition: NGT >500 mL/day lasting beyond post-operative day 3; management: prokinetics, NGT decompression, enteral nutrition via surgical jejunostomy tube placed at the time of pancreaticoduodenectomy or by nasojejunal tube; most cases resolve spontaneously within 2-4 weeks; persistent DGE may require endoscopic dilation or stenting if there is mechanical obstruction at the gastrojejunostomy), post-pancreatectomy hemorrhage (PPH; early <24 hours, late >24 hours; ISGPS classification: grade A no intervention, grade B intervention without ICU, grade C ICU ± reoperation ± embolisation; management: resuscitation, endoscopy or angiography with embolisation, reoperation), and pancreatic ductal adenocarcinoma (PDAC) recurrence (locoregional, liver, peritoneal, distant). Adjuvant chemotherapy: FOLFIRINOX (modified) for 6 months after resection (PRODIGE 24 trial - median survival 54.4 months vs 35.0 months for gemcitabine) or gemcitabine + capecitabine (ESPAC-4 trial). Median survival for resected PDAC is 24-30 months; 5-year survival 15-25\% (improving with modern multi-agent chemotherapy). Borderline resectable PDAC: neoadjuvant chemotherapy (FOLFIRINOX or gemcitabine + nab-paclitaxel) for 2-6 months, then restaging CT and CA 19-9; proceed to surgery if no progression and venous involvement is reconstructable. Arterial involvement (celiac axis, superior mesenteric artery, common hepatic artery) generally precludes resection unless a modified approach (e.g., Appleby procedure for celiac axis involvement) is considered in selected cases.

\medterm{Paraphimosis}-- A urological emergency in which the retracted foreskin cannot be returned to its normal position over the glans penis, causing constriction at the coronal sulcus. The trapped foreskin acts as a tourniquet, causing venous congestion, edema, pain, and eventually arterial compromise and ischemia of the glans. Causes: iatrogenic (failure to reduce the foreskin after catheterisation, cystoscopy, or penile examination—most common), poor hygiene, and balanitis. Predisposing factors: phimosis (narrow, non-retractile foreskin), lichen sclerosus (balanitis xerotica obliterans). Presentation: a painful, swollen glans with a constricting band of retracted foreskin behind the corona. Treatment: (1) Gentle manual reduction—firm, sustained compression of the oedematous glans for 5--10 minutes (ice pack or granulated sugar to reduce edema), then push the glans posteriorly while pulling the foreskin forward. (2) If manual reduction fails—dorsal slit (incision of the constricting band under local anesthesia) or emergency circumcision. (3) Definitive: elective circumcision after resolution of acute inflammation. Complications of paraphimosis: ulceration, necrosis of the glans, and urethral injury.

\medterm{Parathyroid Surgery}-- Primary hyperparathyroidism (PHPT): single adenoma (85\%), hyperplasia (15\%), double adenoma (2\%), carcinoma (<1\%). Criteria for surgery (NIH): age <50, calcium >1 above upper limit, creatinine clearance <60, BMD T-score <-2.5, vertebral fracture, nephrolithiasis, nephrocalcinosis. Localization: sestamibi scan (sensitivity 80-90\%), US (sensitivity 70-80\%), 4D-CT, MRI, PET-Choline. Minimally invasive parathyroidectomy (MIP) for single adenoma with concordant imaging; bilateral exploration for negative or discordant imaging. Intraoperative PTH monitoring (IOPTH): drop of >50\% from baseline 10 minutes after gland excision confirms cure with 95\% positive predictive value. Four-gland exploration: identify 4 glands, resect enlarged, biopsy normal. Auto-transplantation: implant parathyroid tissue (typically 10-15 fragments of 1-2mm each) into brachioradialis or sternocleidomastoid muscle if total parathyroidectomy is necessary (reoperative surgery, secondary/tertiary hyperparathyroidism, multiple gland hyperplasia, MEN1, familial hypocalciuric hypercalcaemia). Reoperative parathyroid surgery: persistent (never normocalcaemic) or recurrent (normocalcaemic for >6 months) PHPT; requires pre-operative imaging (4D-CT, sestamibi with SPECT/CT, MRI, PET-Choline) and intraoperative nerve monitoring (recurrent laryngeal nerve is at risk of injury in reoperative neck surgery). Complications: recurrent laryngeal nerve injury (<1\% in primary surgery, 3-5\% in reoperative), hypoparathyroidism (transient 10-20\%, permanent <2\%), haematoma (1-2\%), infection, and hungry bone syndrome (severe, prolonged hypocalcaemia after resection of severe PHPT with high pre-operative calcium and ALP; management: IV calcium gluconate infusion, oral calcium carbonate 3-6 g/day, calcitriol 0.5-2 mcg/day, magnesium replacement if indicated; worst in patients with pre-operative calcium >3.0 mmol/L and markedly elevated PTH and alkaline phosphatase). Secondary hyperparathyroidism (renal failure): medical management (phosphate binders, vitamin D analogues, cinacalcet); subtotal parathyroidectomy (3.5 gland) or total parathyroidectomy with auto-transplantation for refractory symptomatic hyperparathyroidism with persistent hypercalcaemia despite maximal medical therapy, or for tertiary hyperparathyroidism after kidney transplantation with persistent hypercalcaemia >6 months post-transplant. Tertiary hyperparathyroidism: autonomous PTH secretion despite normal renal function after renal transplantation; parathyroidectomy indicated if hypercalcaemia persists beyond 6 months from transplantation.

\medterm{Penetrating Abdominal Trauma}-- Trauma to the abdomen caused by a missile (gunshot wound) or sharp object (stab wound). GSWs: high energy transfer, produce blast cavity, high risk of intra-abdominal injury (90\%), mandatory exploration. Stab wounds: lower energy, 30-50\% do not penetrate peritoneum; local wound exploration to determine peritoneal violation. Zones: anterior abdomen (xiphoid to pubis, costal margins), flank (midaxillary to posterior axillary line), and back (posterior to posterior axillary line). Diagnosis: CT with triple contrast (IV, oral, rectal) in stable patients; diagnostic laparoscopy for equivocal cases. Management: hemodynamic instability or peritonitis → laparotomy. Stable: serial exams, CT, and selective non-operative management.

\medterm{Phimosis}-- The inability to retract the prepuce (foreskin) over the glans penis. Physiologic phimosis: normal in prepubertal boys, resolves with age. Pathologic phimosis: scarring of the foreskin due to balanitis xerotica obliterans (lichen sclerosus), recurrent infection, or trauma. Presents with ballooning of the foreskin during urination, recurrent balanitis, and urinary tract infections. Complications: paraphimosis (retracted foreskin cannot be reduced, causing glandular ischemia), and in severe cases, urinary obstruction. Treatment: topical corticosteroids (betamethasone) for mild cases; preputioplasty or circumcision for refractory pathologic phimosis. Physiologic phimosis in children does not require treatment.

\medterm{PICC Line (Peripherally Inserted Central Catheter)}-- A long, flexible catheter inserted into the basilic or cephalic vein and advanced so the tip sits in the superior vena cava. Inserted at the bedside, often by a specialized PICC nurse, with ultrasound and ECG or X-ray confirmation. Indications: long-term IV antibiotics (>2 weeks), TPN, chemotherapy, and frequent blood draws. Sizes: 4-6 Fr single, double, or triple lumen. Duration: weeks to months. Advantages over tunneled central lines: lower risk of pneumothorax, insertion by non-physicians. Complications: infection (line-related bloodstream infection), thrombosis, phlebitis, malposition, and catheter fracture.

\medterm{Postoperative Complications}-- Clavien-Dindo classification: Grade I (any deviation, no treatment), Grade II (pharmacologic), Grade III (surgical/endoscopic/radiologic intervention: IIIa no GA, IIIb GA), Grade IV (ICU: IVa single organ, IVb multi-organ), Grade V (death). Common: ileus (delayed GI motility, NPO, NGT, ambulation, gum chewing, neostigmine for Ogilvie), anastomotic leak (clinical: pain, fever, leukocytosis; CT: extraluminal contrast, fluid; management: antibiotics, percutaneous drainage, diversion, reoperation), wound infection (SSI, superficial/deep/organ space; management: antibiotics, drainage, wound care, negative pressure wound therapy or vacuum-assisted closure), surgical site infection prevention bundle, pulmonary complications (pneumonia, atelectasis: IS/CPT, O2, incentive spirometry 10 breaths per hour, early mobilisation), DVT/PE (PESI score, Caprini risk assessment, pharmacologic prophylaxis: enoxaparin 40 mg SC daily or 30 mg BID or unfractionated heparin 5,000 units TID; mechanical prophylaxis: SCD, TED stockings; systemic anticoagulation if PE confirmed). ERAS (Enhanced Recovery After Surgery) pathways: multidisciplinary approach to reduce surgical stress and accelerate recovery; components include pre-operative counseling, avoidance of prolonged fasting (carbohydrate loading 2 hours before surgery), no bowel preparation, epidural analgesia, goal-directed fluid therapy (avoiding excessive crystalloid administration), avoidance of NG tubes and drains, early mobilisation (day of surgery for ambulatory patients, or as soon as the patient can sit in a chair), early oral intake (clear liquids within 4-6 hours, progression to full fluids and then solids as tolerated), and multimodal non-opioid analgesia (paracetamol, NSAIDs, gabapentinoids, lidocaine patches, regional anesthesia blocks); ERAS compliance >70\% associated with reduced length of stay (LOS) and complication rates.

\medterm{Postoperative Ileus}-- Temporary impairment of gastrointestinal motility after surgery, most common after abdominal surgery. Pathophysiology: autonomic dysfunction (sympathetic overactivity), inflammatory mediators, opioids, and electrolyte abnormalities. Presents with abdominal distension, nausea/vomiting, inability to tolerate oral intake, and absence of flatus/stool. Normal postoperative GI recovery: small bowel (24h), stomach (48h), colon (72-120h). Distinguish from mechanical obstruction (pain is crampy, bowel sounds high-pitched). Management: NPO, IV fluids, electrolyte correction, early ambulation, minimize opioids (use multimodal analgesia), gum chewing, and alvimopan (mu-opioid receptor antagonist). Neostigmine for Ogilvie syndrome (colonic pseudo-obstruction).

\medterm{Prostatectomy}-- Surgical removal of the prostate. Radical prostatectomy: for localized prostate cancer (removal of prostate, seminal vesicles, vas deferens, pelvic lymph node dissection). Approaches: open retropubic (RRP), laparoscopic, and robotic-assisted laparoscopic (RALP, most common in the US). Nerve-sparing technique: preserves neurovascular bundles for erectile function and urinary continence. Indications: localized prostate cancer (Gleason \(\ge\)7, PSA >10, or higher stage). Transurethral resection of the prostate (TURP): for BPH (removes central prostate tissue via resectoscope through the urethra), not a cancer operation.

\medterm{Ranula}-- A mucocele (mucus extravasation cyst) of the sublingual salivary gland, presenting as a blue-translucent, fluctuant swelling in the floor of the mouth, typically lateral to the lingual frenulum. A ranula results from trauma or obstruction of a sublingual gland duct, causing mucus accumulation in the sublingual space. If the ranula extends through the mylohyoid muscle into the submental or submandibular space, it is termed a 'plunging' or 'cervical' ranula (presents as a neck mass). Ranulas are benign and painless unless secondarily infected. Differential diagnosis: dermoid cyst, thyroglossal duct cyst, lymphatic malformation, and submandibular gland swelling. Treatment: excision of the ranula and the ipsilateral sublingual gland (lowest recurrence rate—<5\%), marsupialisation (incision and drainage, suturing the cyst wall to the oral mucosa—higher recurrence rate), or laser ablation. Plunging ranulas require transcervical excision of the sublingual gland. Needle aspiration is not curative. Simple oral ranulas can be managed conservatively if small and asymptomatic.

\medterm{Renal Transplant (Renal Transplantation)}-- See Kidney Transplantation.

\medterm{Skin Graft}-- Transfer of skin from a donor site to a recipient site to cover a wound. Split-thickness skin graft (STSG): epidermis + partial dermis (0.008-0.018 inches), harvested with dermatome, donor site heals spontaneously (2-3 weeks). Full-thickness skin graft (FTSG): epidermis + full dermis, donor site requires closure, better cosmetic outcome. Types: autograft (same person), allograft (cadaveric), xenograft (animal), and synthetic (Integra, Biobrane). Uses: burn wounds, traumatic skin loss, chronic ulcers, and post-oncologic reconstruction. Graft take: requires vascularized wound bed, immobilization, and prevention of infection/shear.

\medterm{Small Bowel Obstruction}-- Mechanical blockage of the intestinal lumen. Causes: adhesions (75\%), hernias (15\%), malignancy, Crohn disease strictures, volvulus, intussusception, gallstone ileus, and bezoars. Clinical: crampy abdominal pain, nausea/vomiting (early in proximal SBO), obstipation, and abdominal distension. Signs include high-pitched bowel sounds, localised tenderness (suggesting ischemia), and peritonitis. Laboratory: leukocytosis, elevated lactate (indicator of strangulation), elevated BUN/creatinine (dehydration), and metabolic alkalosis. Imaging: abdominal X-ray shows dilated small bowel loops >3 cm with air-fluid levels. CT abdomen with IV contrast is the gold standard, demonstrating a transition point with proximal dilation and distal collapse, closed-loop obstruction (U-shaped configuration with convergent mesenteric vessels, high risk of ischemia), pneumatosis intestinalis, and signs of perforation; sensitivity >90\% for high-grade SBO. Management: NPO, nasogastric decompression, IV fluids, and CT with contrast. Non-operative observation is appropriate for adhesive SBO without strangulation features (60-80\% resolve over 2-5 days). A water-soluble contrast challenge (Gastrografin) may be both diagnostic (predicts resolution if contrast reaches colon within 4-24 hours) and therapeutic. Operative indications: complete obstruction, peritonitis, strangulation, closed-loop obstruction, failure of non-operative management >48-72 hours, or suspected ischemia. Surgery is via laparoscopy (adhesive SBO) or laparotomy (complicated cases), with adhesiolysis, bowel resection for non-viable segments, and primary anastomosis. Hartmann procedure is reserved for peritonitis or hemodynamic instability. Mortality: 3-5\% for simple obstruction, 25-30\% for strangulated obstruction.

\medterm{Smoke Inhalation Injury}-- Pulmonary injury caused by inhalation of toxic products of combustion. Three components: thermal injury (supraglottic: heat causes edema of the upper airway), chemical irritation (lower airway: aldehydes, ammonia, chlorine, hydrogen chloride cause bronchospasm, mucosal edema, and sloughing), and systemic toxicity (carbon monoxide: carboxyhemoglobin >10\% causes tissue hypoxia; cyanide: impairs cellular respiration). Presents with facial burns, singed nasal hairs, carbonaceous sputum, hoarseness, stridor, and dyspnea. Diagnosis: bronchoscopy (gold standard), carboxyhemoglobin level. Treatment: 100\% oxygen (shortens CO half-life to 60 min); intubation for upper airway edema; mechanical ventilation with lung-protective strategy; bronchodilators; aggressive pulmonary toilet. Hyperbaric oxygen for severe CO poisoning (neuropsychiatric). Hydroxocobalamin for cyanide toxicity.

\medterm{Spinal Fusion}-- Surgical union of two or more vertebrae using bone graft or bone graft substitutes with instrumentation (rods, screws, cages). Approaches: posterior (posterolateral fusion, transforaminal lumbar interbody fusion TLIF), anterior (anterior lumbar interbody fusion ALIF), and lateral (extreme lateral interbody fusion XLIF). Indications: degenerative disc disease, spondylolisthesis, spinal instability, scoliosis, and trauma. Bone graft sources: autograft (iliac crest, gold standard), allograft, BMP (bone morphogenetic protein), and synthetic. Complications: pseudoarthrosis (non-union), hardware failure, adjacent segment disease, infection, and nerve root injury.

\medterm{Splenectomy}-- Indications: trauma (grade IV/V, hemodynamic instability), ITP (refractory), hereditary spherocytosis/elliptocytosis, thalassemia major, sickle cell (sequestration crisis), TTP (refractory), lymphoma staging (historical), hypersplenism, splenic cysts/tumors, Wandering spleen. Laparoscopic preferred. Vaccinations: pneumococcal (PCV20 or PCV15+PPSV23), meningococcal (MenACWY + MenB), Hib - ideally 2 weeks pre-op; post-op if emergency. Antibiotic prophylaxis: penicillin V daily (children <5 or <2 years post-splenectomy or with continued immunosuppression) or amoxicillin daily (adults, alternatively: penicillin V 250 mg BID or amoxicillin 500 mg daily) for at least 2 years after splenectomy or for life if ongoing immunosuppression; immunocompromised patients require lifelong prophylaxis. Fever: stationary IV ceftriaxone plus azithromycin; seek immediate medical attention (splenectomy alert card/wallet card, medical bracelet, fever algorithm for carers). Overwhelming post-splenectomy infection (OPSI): encapsulated bacteria (S. pneumoniae 50-60\%, H. influenzae type b 10-15\%, N. meningitidis 10-15\%, E. coli, Group B Streptococcus, Capnocytophaga canimorsus from dog bites); rapidly progressive (fulminant) sepsis (mortality 50-70\%); presentation: fever, rigors, hypotension, DIC, multi-organ failure within hours; management: IV fluids, broad-spectrum IV antibiotics (ceftriaxone 2 g IV 12-hourly or meropenem 1 g IV 8-hourly) immediately, ICU, vasopressors. Platelet count: rises after splenectomy (thrombocytosis usually transient, peaks at 1-3 weeks, resolves by 2-3 months; platelet count >450,000-1,000,000); aspirin 81 mg daily if platelet count >1,000,000 (to reduce risk of thrombosis). Complications: portal vein thrombosis (3-10\%, increased if splenomegaly, myeloproliferative disorder, or hemolytic anemia; duplex ultrasound at day 7 post-operatively in high-risk patients, or CT with portal venous phase imaging if symptomatic; treatment: therapeutic anticoagulation with LMWH or warfarin for 3-6 months or lifelong if there is thrombophilia or persistent splanchnic vein thrombus). Long-term: re-immunise with PPSV23 every 5 years (or per guidelines; Centers for Disease Control and Prevention recommends a single PPSV23 booster 5 years after the initial dose in asplenic adults, and a second booster at age 65 or above if at least 5 years since last). The risk of OPSI is 0.23-4.4\% per year, with a lifetime risk of 5-10\%.

\medterm{Stoma Complications}-- Complications related to intestinal or urinary stomas. Early (<30 days): ischemia/necrosis (assess with test tube/endoscope, if → mucosa or at fascia → revision), mucocutaneous separation, retraction, parastomal abscess, high output (>1500 mL/day, may cause dehydration/electrolyte depletion → antidiarrheals, loperamide, codeine, octreotide, dietary modification). Late: parastomal hernia (most common late complication, abdominal wall defect → support belt, surgical mesh repair), prolapse (full-thickness protrusion → reduction, repair), stenosis (fibrous ring at skin level → dilation, revision), and skin irritation (leakage → proper appliance fit, skin barriers, powder, paste). Patient education and stoma nurse consultation essential.

\medterm{Surgical Site Infection}-- SSI classification (CDC): superficial incisional (skin/subcutaneous, <30 days), deep incisional (fascia/muscle, <30/90 days with implant), organ/space (any part opened/manipulated, <30/90 days with implant). Risk factors: patient (diabetes, obesity, smoking, immunosuppression, malnutrition, age), operative (contamination class, duration, technique, hair removal, hypothermia, hyperglycemia), postoperative (hyperglycemia, hypothermia, hematoma). Prevention: preoperative shower (CHG), appropriate antibiotic (cefazolin or cefoxitin, or ceftriaxone + metronidazole for colorectal surgery; vancomycin or clindamycin for penicillin-allergic patients) within 60 minutes before incision (120 minutes for vancomycin and fluoroquinolones due to longer infusion time), hair removal with clippers (not razors), normothermia (active warming with forced-air warming blankets), oxygenation (FiO2 80\% during and 6 hours after surgery), maintain euglycaemia (glucose <10 mmol/L, tighter range of 6-10 mmol/L, with insulin sliding scale if needed), wound protector, glove change before closure, and post-operative normothermia, euglycaemia, and oxygen. Antibiotic re-dosing for prolonged surgery (>2 half-lives, approximately 2-4 hours for cefazolin, 4-6 hours for metronidazole, 6-8 hours for ciprofloxacin), appropriate weight-based dosing (e.g., vancomycin 15 mg/kg obese, 30 mg/kg for cefazolin in morbid obesity). Surgical technique: gentle tissue handling, meticulous hemostasis, avoidance of dead space, tension-free closure, use of knotless barbed sutures (e.g., Stratafix, V-Loc) for fascial closure and skin closure techniques: subcuticular absorbable monofilament sutures (skin closure in layers reduces risk of wound infection) or metallic staples (which have the advantage of speed but are not recommended for closure of laparoscopic port sites >5 mm). When closing the fascia, use a continuous slowly absorbable monofilament suture such as PDS (polydioxanone) \#1 or \#2 for laparotomy incisions, a suture length to wound length ratio of at least 4:1 (the "small bites" technique uses 5 mm bites 5 mm apart, reducing incisional hernia). Prophylactic negative pressure wound therapy (ciNPT, Prevena or PICO system) for high-risk wounds (contaminated, dirty, emergency laparotomy, obese, diabetic, stoma, or immunocompromised patients). Management of established SSI: wound culture, open wound, drainage, debridement, packing with saline-soaked gauze or alginate dressings, VAC or negative pressure wound therapy (NPWT) for large, deep, or chronic wounds. Antibiotics only if cellulitis, systemic signs, or organ/space infection. SSI rates: clean <2\%, clean-contaminated <5\%, contaminated <15\%, dirty >25\%.

\medterm{Suture}-- A thread or strand of material used to approximate tissues (bring wound edges together) and hold them in place until healing occurs. Sutures are classified by: (1) Absorbability—absorbable (catgut, polyglactin/Vicryl, poliglecaprone/Monocryl, polydioxanone/PDS—degrade by hydrolysis over weeks to months, used for deep tissues, internal wounds) and non-absorbable (nylon/Ethilon, polypropylene/Prolene, silk, stainless steel—remain permanently, used for skin closure, vascular anastomoses, tendons). (2) Structure—monofilament (single strand—less tissue trauma, lower infection risk, but more difficult to handle, greater knot slippage) and multifilament (braided—easier handling, better knot security, higher infection risk—capillary action harbours bacteria). (3) Needle type—cutting (for skin, tough tissues), reverse cutting, round-bodied (for soft tissues—bowel, vessels, fascia), and taper-cut. Suture size: from 10-0 (ophthalmic microsurgery) to 0--2 (fascia, orthopedics). Suture techniques: simple interrupted (most common, secure, allows individual knot removal), continuous/running (faster, better wound edge eversion, but if one stitch breaks, the entire suture line may fail), horizontal mattress (everts wound edges, tension reduction), vertical mattress (everts edges, good for deep closure), subcuticular (intradermal—best cosmetic result, eliminates suture marks), and purse-string. Alternatives to sutures: surgical staplers (fast, consistent—bowel anastomoses, skin closure), tissue adhesives (cyanoacrylate/Dermabond—for superficial wounds, topical use), adhesive strips (Steri-Strips—for superficial, low-tension wounds), and knotless barbed sutures (Knotless Tissue Control—used in laparoscopic and robotic surgery).

\medterm{TACO}-- Transfusion-Associated Circulatory Overload, a complication of blood transfusion due to volume overload causing pulmonary edema. Risk factors: pre-existing cardiac or renal disease, extremes of age, large transfusion volume, rapid infusion rate. Presents with dyspnea, orthopnea, hypertension, JVD, pulmonary edema, and BNP elevation within 6-12 hours of transfusion. Distinguish from TRALI: TRALI has hypotension, fever, and no fluid overload signs. Management: diuresis (furosemide), oxygen, and slow or split transfusion units in high-risk patients. Prevention: slow transfusion rate, pre-transfusion furosemide.

\medterm{Testicular Torsion}-- A surgical emergency caused by twisting of the spermatic cord, leading to occlusion of the testicular artery and venous outflow, resulting in testicular ischemia and infarction if not treated promptly. Peak incidence: neonatal period and adolescence (12--18 years). Anatomy: the bell-clapper deformity (high insertion of the tunica vaginalis, allowing the testis to rotate freely within the scrotum) is a bilateral predisposing factor present in 12\% of males. Torsion is typically intravaginal (twisting of the testis within the tunica vaginalis). Clinical presentation: sudden onset of severe, unilateral scrotal pain (often waking the patient from sleep), nausea and vomiting, lower abdominal pain, and the affected testis is swollen, tender, and elevated (high-riding testis) with a horizontal lie. The cremasteric reflex is absent on the affected side. Prehn sign (elevation relieves pain in epididymitis) is unreliable. Differential diagnosis: epididymo-orchitis (gradual onset, dysuria, fever, intact cremasteric reflex), torsion of the appendix testis (tender nodule at the superior pole, 'blue dot sign'), and strangulated inguinal hernia. Diagnosis: Doppler ultrasound (reduced or absent testicular blood flow) is the gold standard imaging; sensitivity >95\%. Management: manual detorsion (rotate the testis outward—'open the book'—medial to lateral; success confirmed by pain relief and return of Doppler flow) should be attempted immediately, followed by emergency surgical exploration (scrotal exploration, detorsion, and orchiopexy: fixation of both testes to the scrotal wall with non-absorbable sutures). If the testis is non-viable (black, no bleeding from a small incision): orchiectomy. Salvage rates: >90\% if surgery within 6 hours, 50\% at 12 hours, <10\% after 24 hours. Bilateral orchiopexy is mandatory due to the high rate of bilateral anatomical predisposition.

\medterm{Thoracotomy}-- A surgical incision into the chest wall to access the thoracic cavity (pleural space, lungs, heart, great vessels, esophagus, diaphragm, mediastinum). Types: (1) Posterolateral thoracotomy (most common)—incision below the tip of the scapula, dividing the latissimus dorsi and serratus anterior, entering the chest through the 4th, 5th, or 6th intercostal space; provides excellent access for pulmonary lobectomy, pneumonectomy, oesophagectomy, and thoracic aortic surgery. (2) Anterolateral thoracotomy—emergency access for trauma, pericardial window, or lung biopsy. (3) Median sternotomy—midline division of the sternum (sternotomy), used for cardiac surgery (CABG, valve surgery, thymectomy, anterior mediastinal mass). (4) Clamshell incision (bilateral anterolateral thoracotomy with transverse sternotomy)—for bilateral lung transplantation, massive trauma. (5) Axillary thoracotomy—for first rib resection, sympathectomy. Complications: post-thoracotomy pain (significant, long-lasting—use of epidural analgesia, paravertebral block, intercostal nerve blocks), wound infection, hemorrhage, persistent air leak, pneumonia, and atelectasis. Minimally invasive alternatives: video-assisted thoracoscopic surgery (VATS)—reduces pain, hospital stay, and recovery time.

\medterm{Thyroid Surgery}-- lobectomy, subtotal thyroidectomy, and total thyroidectomy. Indications: Bethesda V/VI nodules on cytology, thyroid cancer (papillary >1.5 cm, follicular, medullary, anaplastic), Graves disease (compressive symptoms, failed medical therapy, severe ophthalmopathy), and multinodular goiter with compressive symptoms or substernal extension. Lobectomy is appropriate for Bethesda III/IV nodules and low-risk papillary microcarcinoma (<1.5 cm, solitary, no nodal involvement). Total thyroidectomy is standard for bilateral disease, high-risk features, medullary carcinoma, and anaplastic carcinoma (palliative for airway obstruction). Central compartment dissection (level VI) is performed for clinically involved nodes; prophylactic dissection is controversial due to increased risk of nerve injury and hypoparathyroidism. Lateral neck dissection (levels II-V) is reserved for confirmed cervical metastases. Intraoperative nerve monitoring facilitates identification of the recurrent laryngeal nerve. The external branch of the superior laryngeal nerve is at risk during ligation of the superior thyroid pole. The parathyroid glands are at risk of devascularization or accidental excision. Complications: recurrent laryngeal nerve injury (2-5\% transient, <1\% permanent) presents with hoarseness and aspiration; bilateral injury causes stridor requiring tracheostomy. Hypoparathyroidism (transient 10-30\%, permanent 1-2\%) requires calcium and calcitriol supplementation, guided by post-operative PTH monitoring. Haematoma (0.5-2\%) requires emergency evacuation to prevent airway compromise. Wound infection and hypertrophic scarring are less common. TSH suppression therapy reduces recurrence risk in thyroid cancer. Radioactive iodine ablation is given for high-risk differentiated cancer after total thyroidectomy.

\medterm{TIPS (Transjugular Intrahepatic Portosystemic Shunt)}-- A radiologically placed shunt connecting the portal vein to a hepatic vein within the liver, creating a portosystemic shunt to reduce portal pressure. Indications: variceal bleeding (refractory or recurrent), refractory ascites, and hepatic hydrothorax. Procedure: transjugular access, hepatic vein cannulation, needle puncture through liver parenchyma into the portal vein, balloon dilation, and covered stent placement. Portal pressure gradient goal: <12 mmHg. Contraindications: severe pulmonary hypertension, right heart failure, severe hepatic encephalopathy, and severe liver failure (MELD >18-20). Complications: encephalopathy (30-40\%), shunt stenosis/occlusion, bleeding, infection, and radiation exposure.

\medterm{Total Hip Arthroplasty (THA)}-- Replacement of the hip joint with prosthetic components. Acetabular component: metal shell + polyethylene liner. Femoral component: metal stem + femoral head (cobalt-chrome or ceramic). Fixation: cemented (PMMA, for elderly, osteoporotic bone), cementless (porous-coated, for younger patients, biologic fixation), and hybrid. Approaches: posterior (most common, dislocation risk 2-5\%), direct anterior (DAA, muscle-sparing, less pain), anterolateral, and lateral (Hardinge). Indications: end-stage osteoarthritis, avascular necrosis, rheumatoid arthritis, and hip fracture (femoral neck). Complications: dislocation, infection (1-2\%), VTE, periprosthetic fracture, leg length discrepancy, and aseptic loosening (long-term).

\medterm{Total Knee Arthroplasty (TKA)}-- Replacement of the knee joint with prosthetic components. Femoral component (cobalt-chrome), tibial component (titanium tray + polyethylene insert), and patellar component (polyethylene, optional). Fixation: cement (PMMA) most common in older adults; cementless in younger. Alignment: mechanical (neutral, most common), kinematic, and anatomic. Indications: end-stage osteoarthritis (most common), inflammatory arthritis, and post-traumatic arthritis. Complications: infection (1-2\%), stiffness (manipulation under anesthesia), instability, patellar complications (maltracking, fracture), periprosthetic fracture, VTE, and aseptic loosening.

\medterm{Tracheostomy}-- A surgical opening (stoma) in the trachea through the neck. Percutaneous dilational tracheostomy (PDT, ICU bedside, Seldinger technique, most common) vs. open surgical tracheostomy (OR, for difficult anatomy, large goiter, obese neck). Indications: prolonged mechanical ventilation (expected >14-21 days, reduces VAP, sedation, and ICU length of stay), upper airway obstruction, and severe secretion management. Tracheostomy tube: cuffed (for ventilation) or uncuffed/fenestrated (for spontaneous breathing and speech). Complications: bleeding, infection, pneumothorax, subcutaneous emphysema, tracheal stenosis, tracheomalacia, and stomal infection.

\medterm{Tracheotomy}-- See Tracheostomy.

\medterm{TRALI}-- Transfusion-Related Acute Lung Injury, a life-threatening complication of blood transfusion characterized by acute hypoxemia and bilateral pulmonary edema within 6 hours of transfusion, without evidence of volume overload. Pathophysiology: donor antibodies (anti-HLA, anti-HNA) in plasma-rich products (FFP, platelets) activate recipient neutrophils, causing endothelial injury and capillary leak in the lungs. Presents with dyspnea, hypoxemia, fever, and bilateral infiltrates on CXR (mimics ARDS). Management: supportive care (oxygen, mechanical ventilation), similar to ARDS. Prevention: use male-only plasma to reduce antibody risk. Mortality 5-10\%.

\medterm{Transfusion Reaction}-- An adverse reaction to transfusion of blood products. Acute hemolytic transfusion reaction: ABO incompatibility, immediate fever/chills, hypotension, back pain, DIC, renal failure; management: STOP transfusion, maintain BP/urine output, supportive care. Febrile non-hemolytic transfusion reaction: most common, fever/chills without hemolysis. Allergic: urticaria; treat with antihistamines. Anaphylactic: severe allergic, IgA deficiency; treat with epinephrine. Transfusion-related acute lung injury (TRALI): acute hypoxemia, bilateral pulmonary edema within 6 hours; supportive care. Transfusion-associated circulatory overload (TACO): pulmonary edema from volume overload; diuretics. Bacterial contamination: sepsis from contaminated unit.

\medterm{Transplant Rejection}-- Immune-mediated destruction of a transplanted organ or tissue, classified by timing and mechanism. Hyperacute: preformed donor-specific antibodies, minutes to hours post-transplant, immediate graft thrombosis, irreversible; prevention by cross-matching. Acute: T-cell mediated (most common, days to weeks/months), responds to immunosuppression; presents with graft dysfunction and biopsy findings. Chronic: antibody-mediated, progressive, irreversible fibrosis and graft dysfunction over months to years. Diagnosis: graft biopsy (Banff classification). Treatment: acute rejection — pulse corticosteroids, thymoglobulin, IVIG, plasmapheresis; chronic — optimize immunosuppression, manage risk factors.

\medterm{Transplant Surgery}-- Transplant surgery involves replacing a failing organ with a healthy donor organ. Kidney transplant: living donors (laparoscopic nephrectomy, left kidney preferred) or deceased donors (standard criteria, expanded criteria, DCD). Induction immunosuppression includes basiliximab (standard risk) or thymoglobulin/alemtuzumab (high immunologic risk). Maintenance typically uses tacrolimus, mycophenolate, and prednisone; belatacept offers CNI-free options for EBV-positive patients. Rejection: hyperacute (preformed antibodies, immediate graft loss), acute cellular (T-cell mediated, responds to steroid pulse or ATG), acute antibody-mediated (C4d+ on biopsy with donor-specific antibodies, treated with IVIG, plasmapheresis, rituximab), and chronic allograft nephropathy (interstitial fibrosis and tubular atrophy). Liver transplant allocation uses the MELD-Na score. Donor hepatectomy may be whole or split/living donor. Recipient hepatectomy uses bicaval or piggyback technique with portal vein, hepatic artery, and biliary reconstruction (duct-to-duct or Roux-en-Y hepaticojejunostomy). Complications: primary non-function, hepatic artery thrombosis (3-5\%, requires retransplantation unless collateralized), biliary leak/stricture, portal vein thrombosis, and rejection. PTLD is an EBV-driven lymphoproliferative complication managed by immunosuppression reduction and rituximab. Pancreas transplant (SPK, PAK, PTA) is performed for type 1 diabetes with ESRD; enteric drainage is preferred for exocrine secretion. Long-term management: immunosuppression optimization, infection prophylaxis (CMV, EBV, fungal), and surveillance for graft dysfunction, PTLD, and cardiovascular disease. One-year graft survival exceeds 90\% for kidney and 85\% for liver; 5-year survival >70\% for both.

\medterm{Trauma Surgery}-- Trauma surgery encompasses the acute care of injured patients, guided by ATLS (ABCDE) principles: Airway maintenance with cervical spine protection, Breathing and ventilation, Circulation with hemorrhage control, Disability/neurologic assessment, and Exposure/environmental control. Hemorrhage control is the priority, achieved via direct pressure, tourniquets for extremity injuries, pelvic binders for pelvic fractures, and REBOA for exsanguinating torso hemorrhage. Damage control surgery is indicated for patients with the lethal triad of hypothermia, acidosis, and coagulopathy. Phase 1 involves abbreviated laparotomy with bleeding control (packing, shunting) and contamination control. Phase 2 is ICU resuscitation to correct hypothermia (active rewarming), coagulopathy (massive transfusion at 1:1:1 ratio of pRBCs, FFP, platelets; tranexamic acid 1 g IV within 3 hours of injury), and acidosis. Phase 3 is return to the OR for definitive repair after physiologic normalization. Damage control resuscitation incorporates permissive hypotension (SBP 80-90 mmHg) in penetrating torso trauma without head injury. Definitive surgical management depends on injury pattern and may include splenectomy, liver packing and resection, bowel resection, vascular repair, and fracture fixation. Resuscitative thoracotomy is indicated for penetrating chest trauma with witnessed arrest and short transport time (best survival for cardiac tamponade). Tension pneumothorax requires immediate needle decompression (14G, 2nd ICS midclavicular line) followed by tube thoracostomy. Massive haemothorax (>1500 mL initial output) requires expeditious thoracotomy. Abdominal compartment syndrome (bladder pressure >20 mmHg with organ dysfunction) may necessitate decompressive laparotomy with temporary closure. FAST examination and CT imaging guide surgical decision-making. Enteral nutrition is initiated within 24-48 hours. Outcome depends on injury severity, promptness of intervention, and successful resuscitation.

\medterm{Traumatic Brain Injury}-- Disruption of brain function caused by an external mechanical force. Severity: mild (GCS 13-15, concussion), moderate (GCS 9-12), severe (GCS 3-8). Primary injury: direct damage at moment of impact (contusion, diffuse axonal injury, intracranial hemorrhage). Secondary injury: evolves over hours to days (cerebral edema, increased ICP, hypoxia, hypotension, ischemia, seizures). Intracranial hemorrhage: epidural (arterial, lenticular), subdural (venous, crescentic), intracerebral, subarachnoid. Management: ABCs, prevent secondary injury (maintain SBP >90, SpO2 >94\%, PaCO2 35-45), ICP monitoring (goal <22), CPP >60 mmHg, head of bed 30°, hyperosmolar therapy (mannitol, hypertonic saline), seizure prophylaxis, and surgical evacuation of mass lesions.

\medterm{Trephining (Trepanation)}-- A surgical procedure in which a hole is drilled or scraped into the skull (cranium), exposing the dura mater. Historically, trephining dates back to Neolithic times (10,000 BCE) and was performed for ritual, therapeutic (releasing 'evil spirits'), and post-traumatic indications. Modern medical trephination: (1) Burr hole—small craniotomy using a handheld or electric drill (burr) for evacuation of intracranial haematomas (epidural, subdural—including twist-drill craniostomy for chronic subdural haematoma), placement of external ventricular drain (EVD) for intracranial pressure monitoring or CSF drainage, and brain biopsy. (2) Decompressive craniectomy—removal of a large skull flap to treat refractory intracranial hypertension (traumatic brain injury, malignant middle cerebral artery infarction, cerebral venous sinus thrombosis). (3) Craniotomy—larger opening for definitive neurosurgery (tumor resection, aneurysm clipping, epilepsy surgery). (4) Stereotactic trephination—frame-based or frameless stereotactic navigation for biopsy, deep brain stimulation electrode placement, or tumor ablation. Instruments: Hudson brace, perforator and burr, craniotome, and high-speed drills. Complications: bleeding (epidural, subdural, intracerebral), infection (osteomyelitis, meningitis, brain abscess), CSF leak, pneumocephalus, seizures, and brain herniation.

\medterm{Tubal Ligation}-- Surgical sterilization by occluding the fallopian tubes, preventing the ovum from reaching the uterus. Postpartum tubal ligation: after vaginal or cesarean delivery (Pomeroy technique, Parkland). Interval: laparoscopic (bands, clips, bipolar cautery) or mini-laparotomy (Hulka or Filshie clip). Reversal possible with tubal reanastomosis (pregnancy rate 50-80\%). Female sterilization: 99\% effective, permanent. Complications: ectopic pregnancy if procedure fails (higher after bipolar cautery than clips). Essure (hysteroscopic coils) removed from the market.

\medterm{Umbilical Hernia}-- A hernia through the umbilical ring, common in infants (most umbilical hernias close spontaneously by age 4-5) and adults (risk factors: obesity, pregnancy, ascites, cirrhosis). Presents as a reducible bulge at the umbilicus, worse with increased intra-abdominal pressure. In adults: risk of incarceration and strangulation is higher. Diagnosis: clinical. Treatment: in children: observe, surgically repair if persists >age 4-5 or >2 cm. In adults: surgical repair (primary suture repair, mesh onlay, or laparoscopic IPOM) due to risk of complications.

\medterm{Ureteral Stent (Double-J Stent)}-- A tube placed in the ureter from the renal pelvis to the bladder to maintain patency, with pigtail coils at both ends (hence double-J). Indications: ureteral obstruction (stones, stricture, malignancy), after ureteroscopy (prevent post-op obstruction), and as a pre-treatment for complex ureteral stones. Placed cystoscopically or antegrade through a nephrostomy. Sizes: 6-8 Fr, lengths 22-30 cm. Left in for weeks to months (though should not exceed 3 months due to encrustation). Complications: stent-related discomfort (urgency, frequency, hematuria, flank pain — very common), infection, migration, encrustation (forgotten stent), and stone formation.

\medterm{Urethral Stricture}-- Narrowing of the urethra due to fibrosis of the corpus spongiosum, most commonly caused by trauma (pelvic fracture, straddle injury, iatrogenic from catheterisation or urethral instrumentation), infection (gonococcal urethritis leading to stricture years later, now less common in the antibiotic era), and lichen sclerosus (Balanitis Xerotica Obliterans, BXO). Presents with progressive urinary obstruction: weak stream, hesitancy, straining, incomplete bladder emptying, frequency, and recurrent UTIs. Complete obstruction causes acute urinary retention. Diagnosis: retrograde urethrogram demonstrates the location and length of the stricture. Urethroscopy confirms and may allow simultaneous treatment. Treatment depends on location, length, and severity: urethral dilation (temporary, high recurrence rate), internal urethrotomy (optical incision of the stricture, 50\% recurrence at 2 years for bulbar strictures <2 cm), and urethroplasty (open surgical reconstruction, the gold standard for recurrent or long strictures, success rates >90\% for primary bulbar repairs). Complications include recurrence, erectile dysfunction, and incontinence.

\medterm{Urology}-- Surgical specialty focused on the diagnosis and treatment of diseases of the male and female urinary tract (kidneys, ureters, bladder, urethra) and the male reproductive system (prostate, penis, testes, seminal vesicles). Subspecialties: urologic oncology (prostate, bladder, kidney, testicular cancer), endourology (minimally invasive stone surgery, ureteroscopy, PCNL), female urology (incontinence, pelvic organ prolapse), neurourology (voiding dysfunction, neurogenic bladder), andrology (male infertility, erectile dysfunction), and pediatric urology (congenital anomalies, hypospadias, undescended testes). Common procedures: transurethral resection of the prostate (TURP), cystoscopy, vasectomy, nephrectomy, and radical prostatectomy.

\medterm{Varicocele}-- Abnormal dilatation of the pampiniform plexus of veins in the spermatic cord, caused by incompetent or absent valves in the internal spermatic (testicular) vein. Left-sided varicoceles are far more common (85-90\%) because the left testicular vein inserts into the left renal vein at a right angle, whereas the right testicular vein inserts obliquely into the inferior vena cava. Acute-onset right-sided varicocele raises suspicion for retroperitoneal pathology (renal cell carcinoma with tumor thrombus extending into the IVC, retroperitoneal mass). Presents as a scrotal swelling described as a bag of worms on palpation, typically painless or causing a dull ache/heavy sensation that worsens with standing or Valsalva and improves with lying supine. Grading: Grade I (palpable only with Valsalva), Grade II (palpable without Valsalva), Grade III (visible through the scrotal skin). Varicocele is the most common correctable cause of male infertility: elevated testicular temperature from venous congestion impairs spermatogenesis (reduced sperm count, motility, and increased abnormal morphology). Diagnosis: clinical examination and scrotal ultrasound with Doppler. Treatment: indicated for symptomatic varicocele, infertility with abnormal semen analysis, testicular hypotrophy in adolescents, or patient preference. Options: microscopic varicocelectomy (gold standard, highest success, lowest recurrence), laparoscopic varicocelectomy, or percutaneous embolization. Subclinical varicoceles (detected only on ultrasound) do not require treatment.

\medterm{Vascular Access}-- Vascular access for hemodialysis includes arteriovenous fistulas (AVF), arteriovenous grafts (AVG), and central venous catheters (CVC). The fistula first initiative prioritizes AVF as the gold standard due to best long-term patency and lowest complication rates. Common configurations: radiocephalic (wrist), brachiocephalic (elbow), and brachiobasilic transposition (upper arm). Preoperative vein mapping by ultrasound is essential for planning. AVGs (PTFE) mature faster (2-4 weeks vs. 2-4 months for AVF) and are useful when native veins are inadequate, but have higher rates of thrombosis and infection. Tunneled CVCs allow immediate use and serve as a bridge to maturation or for patients with exhausted peripheral options, but carry the highest infection and central venous stenosis risk. Non-tunneled CVCs are for short-term emergent use. Maturation criteria: vein diameter >6 mm, depth <6 mm, flow >600 mL/min. Indications: end-stage renal disease requiring hemodialysis. Complications: stenosis (angioplasty, stenting), thrombosis (thrombectomy, catheter-directed thrombolysis), steal syndrome (hand ischemia managed by DRIL or PAI procedures), aneurysm and pseudoaneurysm (resection, ligation), infection (graft excision, antibiotics), and high-output cardiac failure (revision, banding, or ligation). CVC complications: pneumothorax (reduced by ultrasound guidance), arterial puncture, catheter-related bloodstream infection (removal, tip culture, antibiotics), central venous stenosis, and air embolism. Patency: AVF > AVG > CVC.

\medterm{Vasectomy}-- Male surgical sterilization: cutting and sealing the vas deferens. No-scalpel vasectomy: puncture technique, less bleeding, faster recovery. Conventional: small incisions. Fascial interposition: prevents recanalization. Post-procedure: semen analysis at 12 weeks or 20 ejaculates (azoospermia confirms sterility). Failure rate: <1\% (from recanalization). Reversal: vasovasostomy (microsurgical, success 70-95\% depending on time since vasectomy). Complications: hematoma, infection, chronic scrotal pain (congestive epididymitis, post-vasectomy pain syndrome, 1-2\%), and sperm granuloma.

\medterm{Venous Stasis Ulcer}-- A chronic wound caused by chronic venous insufficiency, most common type of leg ulcer (70-80\%). Located on the medial gaiter area (medial malleolus). Presents with shallow, irregular ulcer with granulation tissue, surrounded by hemosiderin deposition (brownish discoloration), lipodermatosclerosis, edema, and varicose veins. Often painful and associated with venous insufficiency (CEAP classification). Diagnosis: clinical; venous duplex ultrasound to assess reflux. Treatment: compression therapy (multi-layer bandaging, graded compression stockings 30-40 mmHg) is the mainstay. Adjunctive: leg elevation, wound debridement, moist wound healing, and pentoxifylline. Refractory: venous ablation or surgery.

\medterm{Venous Ulcer}-- See Venous Stasis Ulcer.

\medterm{Wound Dehiscence}-- Separation of the layers of a surgical wound after closure, most commonly occurring within 3-14 days postoperatively. Risk factors: wound infection, poor nutrition (low albumin), immunosuppression, obesity, diabetes, smoking, increased intra-abdominal pressure (coughing, straining, ileus), and poor surgical technique (tight sutures, extensive electrocautery). Presents with sudden serosanguinous drainage from the wound, often with a pop sensation. Superficial dehiscence: skin and subcutaneous layers open. Deep (fascial) dehiscence: separation of the fascia, often with evisceration (viscera protruding through the wound) → surgical emergency. Management: superficial → wound care and NPWT. Deep → emergent return to OR for closure with retention sutures.

\medterm{Xenograft (Heterograft)}-- A surgical graft taken from a donor of a different species (xenos = foreign, Greek). Xenografts are used as temporary biological dressings or as tissue replacements. The most common xenograft source is porcine (pig) skin, used as a temporary wound covering for extensive burns (partial-thickness and deep partial-thickness burns) and large traumatic skin loss. Porcine xenografts adhere to the wound bed, reduce fluid and protein loss, decrease pain, and provide a barrier against bacterial contamination while the patient's own skin regenerates or until autografting is possible. Other xenograft sources include bovine (cow) pericardium (used in cardiac surgery for valve reconstruction, pericardial patch repair, and as a bioprosthetic heart valve material — Carpentier-Edwards, Hancock valves), and bovine bone grafts (used in orthopaedic and dental surgery for bone void filling, spinal fusion, and maxillofacial reconstruction). Xenografts trigger a strong immune response (hyperacute rejection mediated by pre-existing xenoreactive antibodies against Gal-alpha-1,3-Gal epitopes on porcine cells) and are therefore rejected within 7-14 days if used as viable tissue, limiting their use to temporary coverage. Genetically modified pigs (knockout of Gal epitopes) have been developed to reduce immunogenicity and enable longer-term xenograft survival. The first pig-to-human heart transplant (xenotransplantation) was performed in 2022 (University of Maryland, 10 gene-modified pig heart, recipient survived 2 months). Compared to allografts (human donor grafts) and autografts (patient's own tissue), xenografts have the advantages of unlimited availability and no donor site morbidity, but are inferior in durability and immunocompatibility for permanent replacement.

